\section{Implications for holography in de Sitter space}\label{sec:implications}

Sections~\ref{sec:methods} and~\ref{sec:derivations} derived AdS/CFT from a single
microstate, with no ensemble average. The present section shows that this apparatus
meets the criteria for holography in de~Sitter space already fixed in the literature.
The aim is not to propose new physics: the mechanisms it uses---the tensors of general
relativity, the Polyakov worldsheet, Kaluza--Klein compactification, the warped throat,
Jacobson's horizon thermodynamics---are established results. What is shown is that these
principles already coincide with the microstate; much of the section is therefore a
review of the known principle followed by how the object realises it.

\subsection{What the section does, and against what it is measured}\label{ssec:noselect}

Holography in de~Sitter, as Witten poses it~\cite{Witten01dS}, requires that general relativity be embedded
in a more precise theory that determines the cosmological constant $\Lambda$, and in
which the observer is described by the theory rather than introduced from
outside---because in a cosmology with $\Lambda>0$ there is no asymptotic region to which
an external observer may retreat.

The yardstick is the classification of von Neumann factors into three types---used in
this context by Chandrasekaran, Longo, Penington and Witten---three modes of existence of
the same field. \emph{Type~I} is that of discrete objects---particle, path integral,
string counted one by one, operator of definite spectrum---with minimal projections.
\emph{Type~II} is that of observation---the trace, the crossed product of the algebra by
its modular flow---where observing turns an algebra without trace into one with a trace.
\emph{Type~III} is that of the local field---the wedge, the horizon, the worldsheet, the
spacetime---without normal trace or projections: non-locality.

Of these types, what the section \emph{proves} is the signature of each, not the
technical classification: type~I as discrete spectrum (\ProofWitness{eigenvalues\_modulus\_half, witten\_finite\_rank}); type~II as an invariant the flow preserves---the trace in
concrete form (\ProofWitness{omega\_flow\_invariant, tower\_scale\_group}); type~III as the
non-factorisation of the horizon into interior${}\otimes{}$exterior
(Proposition~\ref{prop:obs-nofirewall}). The technical names (II$_1$, III$_1$) and the
crossed product are taken from CLPW by citation; the signatures are proved here.
Signature first, name second.

The criteria come from three distinct sources, which should not be conflated. From
\textbf{Witten 2001} (\emph{Quantum Gravity in De Sitter Space}, W01) come the
determination of $\Lambda$, the finiteness of the Hilbert space with $S=\ln N$, the
isometry-group obstruction, the pairing $\langle f|i\rangle$, the CPT map, the two
conjectures---of unitarity and of entropy---and the antipodal singularity. From
\textbf{Chandrasekaran, Longo, Penington and Witten 2022} (CLPW) come the observer
described by the theory, the maximally mixed vacuum state, and the construction of the
type~II algebra by crossed product from the type~III. The thermal nature of the static
patch predates both: it is the result of \textbf{Gibbons--Hawking}. Throughout, each
criterion is tagged with which of the three sources states it.

The thesis is that the framework realises the three types on a single object, the
microstate $\Om$ of modulus $|\Om|=\tfrac12$, and thereby satisfies the requirements
Witten left open. The three types are not a label per subsection but three threads
crossing the whole section: type~I runs through the particle (\S\ref{ssec:hinge}), its
path integral, and the discrete string of the tower (\S\ref{ssec:accum})---to it belongs
the arithmetic operator of the zeros, referred to the Hilbert--P\'olya companion and not proved
here (the tower's own Hamiltonian, self-adjoint and with a gap, is
Proposition~\ref{prop:operators} and Theorem~\ref{thm:intertwine}); type~II runs
through the ER${=}$EPR bridge (\S\ref{ssec:hinge}), observation
(\S\ref{ssec:observer}) and the tracial flow (\S\ref{ssec:accum}); type~III runs through
the worldsheet (\S\ref{ssec:hinge}) and the spacetime with gravity (\S\ref{ssec:accum}).
The modulus $\tfrac12$, the projection $\Pi$, the entropy $S$, the centre $C$ and the
conjugation $\Theta$ recur across movements; each recurrence declares its register and
its anchor. The repetition is the weave of a multi-thread demonstration.


\subsection{The non-local field and the particle that emerges from it}\label{ssec:hinge}

\paragraph{The worldsheet, the type~III field, and its link to gravity.}
General relativity requires that the laws not depend on coordinates; the tensor guarantees
this, because a tensor equation true in one frame is true in all. In string theory this
invariance lives on the worldsheet: the string sweeps a surface whose dynamics is the
Polyakov action, written on the worldsheet metric $h_{ab}$---a tensor---and invariant under
diffeomorphisms and Weyl rescalings; its constant is the tension $T=1/2\pi\alpha'$. The
worldsheet is the type~III field of the section: a local wedge without normal trace, and it
is the same field that will reappear as the curved spacetime of \S\ref{ssec:accum}---the
type~III face that opens and closes \S\ref{sec:implications}. Its non-locality is geometric,
not postulated. The worldline $\Om(\tau)=\tfrac12 e^{i\tau\ln\golden}$, promoted by Polyakov,
closes onto the torus $T^2_{\PCF}=\C/\Lambda$ with modular parameter $\tau=i$
fixed---the self-dual point of S-duality, not chosen---and $\golden$ generates the lattice.
A torus is compact and boundaryless, and a boundaryless domain admits no question of what
happens at the edge of one region against another, because there is no edge at which to
separate (\ProofWitness{fibre\_independent\_of\_base}). By the modulus $\tfrac12$, the loop
around the torus accumulates a half rotation, $e^{i\pi}=-1$: the sheet is a M\"obius band,
non-orientable, with no privileged interior side. The particle is the bulk--boundary
interface this field carries:
\begin{equation}\label{eq:obs-interface}
  \Pi\colon\E^{3}\longrightarrow\C ,
\end{equation}
whose Hausdorff dimension is \eqref{eq:hausdorff}, established in \S\ref{ssec:construction}.
This $\Pi$ is the projection of the first movement: in closed form
$\pi_{\PCF}(a,b,c)=(ab/c)\,\pi/(3\sqrt3)$ \eqref{eq:projpcf}, the one constant that fixed
$\varepsilon_0$ \eqref{eq:eps0} and the Fibonacci sum (Proposition~\ref{prop:oplus}).
\ProofWitnessBlock{projection\_closed\_form, epsilon0\_from\_projection}%
The horizon of $\Pi$ does not factorise into interior${}\otimes{}$exterior, which is the
type~III signature and the reason there is no firewall:
\begin{proposition}[Holographic--relativistic horizon (no firewall)]\label{prop:obs-nofirewall}
The projection $\Pi\colon\E^{3}\twoheadrightarrow\C$ creates a horizon that is at
once \emph{holographic} (information fully encoded on the boundary) and
\emph{relativistic} (the Galois conjugation $\golden\leftrightarrow\bar\golden$
exchanges the two sides like the two wedges of a Rindler spacetime): $\Pi$ is
surjective, with $\ker\Pi=\{(0,0,0)\}$ and $\Pi^{-1}(p)\neq\varnothing$ for every
$p\in\C$. No averaging over microstates is needed: the microstate is one state and not an ensemble, as established in \S\ref{ssec:adscft} (its realisation as a self-adjoint operator being \S\ref{sec:dof})---so the AMPS firewall ambiguity~\cite{AMPS13} does not arise in this setting: with a single microstate there is no ensemble of typical states for the argument to average over (\S\ref{sssec:ensemble}).
\ProofWitnessBlock{omega\_flow\_invariant, bridge\_cocycle, bridge\_inverse, bridge\_compose}
\end{proposition}
Three independent facts confirm the non-separability, and all three fall on $\tfrac12$. At
every scale, $|\Om|_\sigma=\tfrac12\ \forall\sigma$ (\ProofWitness{omega\_flow\_invariant}). At
every separation, the entanglement is $|\Om(\tau_A)\bar\Om(\tau_B)|=\mu_3^2=\tfrac14$, and
this ER${=}$EPR bridge is derived, not postulated: the operators $T(\sigma_1,\sigma_2)$ form
a group (Proposition~\ref{prop:er-epr};
\ProofWitnessBlock{bridge\_cocycle, bridge\_inverse, bridge\_compose}).

The conjugation $\Theta={}$CPT, the antilinear Galois map $\golden\leftrightarrow\bar\golden$,
swaps the wedges $P\leftrightarrow F$ and fixes $C$.

\emph{W6 --- the antilinear $\Theta$/CPT map (W01).} Witten defines the Hilbert space
non-perturbatively through an antilinear map $\Theta$ (a CPT operation) that renders the
bilinear pairing Hermitian (W01, eq.~18). It is the Galois conjugation
$\golden\leftrightarrow\bar\golden$: $\Theta$ swaps the two Rindler wedges and fixes the
centre, with $\Theta^2=1$; it preserves the modulus, $\lVert\Theta z\rVert=\lVert z\rVert$, hence is antiunitary. \emph{Proof:}
\ProofWitnessBlock{cpt\_swaps\_P\_to\_F, cpt\_fixes\_C, cpt\_involution, cpt\_preserves\_modulus}.

\emph{W7 --- the $\mathrm{SO}(1,n{-}1)$ obstruction (W01).} With $N$ finite the de~Sitter
isometry group cannot act on the Hilbert space; being non-compact it has no
finite-dimensional unitary representations (W01). Its place is taken by the discrete
$\mathbb{Z}_3$ symmetry of the cube roots of unity with the scale invariance
$|\Om|_\sigma=\tfrac12$. \cite{Witten01dS}. \ProofWitnessBlock{eisenstein\_cube, omega\_conj}

\paragraph{The whole geometry as a tensor: the warp generates curvature.}
This spacetime is tensorial in its entirety. The throat that compactifies the scale has a
warped metric with radial
coordinate $z(\sigma)=\golden^{\sigma}$ and non-vanishing warp gradient: with $A=-\log z$, so
$A'=-1$, the standard machinery---from $g$ to Christoffel, to Riemann, to Ricci---turns that
gradient into curvature computed, not assumed:
$R_{\mu\nu}=-(A''+d\,A'^2)\,g_{\mu\nu}=-4\,g_{\mu\nu}$, so the throat is an Einstein space
(\S\ref{app:curvature}; \ProofWitness{R\_scalar\_pcf, R\_AB\_einstein}). The densification of
the warp \emph{is} the Ricci tensor. The extrinsic side closes it: the constant-$\sigma$
leaves, their embedding and second fundamental form are tied to the intrinsic curvature by
Gauss's equation---de~Sitter is the hyperboloid whose intrinsic curvature is the wedge of
the extrinsic curvature of the embedding, $R_{\mu\nu}=3H^2 g_{\mu\nu}$, developed in full in
\S\ref{app:embedding} (\ProofWitness{ds\_ricci\_from\_gauss}).

\paragraph{From that same geometry the particle emerges, discrete (type~I).}
The warp is sustained by a stress--energy tensor: Einstein's equation with source balances
the throat geometry against a matter tensor whose spatial part is pressure. The tensor that
in gravity describes the metric is the same that describes the tension---both words share
the root \emph{tendere}---and it is that tension that sustains the throat. Geometrically the
throat is a densifying centre: on the Riemann sphere, with zero at one pole and infinity at
the other, the golden tower produces concentric shells growing outward and a centre
compressed without bottom, since $e^{2A}=z^{-2}=\golden^{-2\sigma}$ diverges as $z\to0$. The
centre is superpoint $\mathbb{R}^{0|1}$-like, the object without metric or spin structure,
with a single odd direction (\S\ref{sssec:superpoint}). Its $\mathbb{Z}_2=\{0,1\}$ grading
distinguishes the even, bosonic directions (the ``0'') from the single odd, fermionic one
(the ``1''); the centre of the tension is that odd direction: fermionic. The involution
$\alpha(x)=1-x$ swaps $0$ and $1$ and has no fixed point between them---which makes the
binary register self-contradictory---but a unique fixed point at $x=\tfrac12$: the extremes
$\{0,1\}$ are the binary register---the bit, the boson, the fermion---and the self-conjugate
midpoint $\tfrac12$ is at once the modulus $|\Om|$ and the fermion's \emph{spin} $\tfrac12$.
The particle is thus coded by three readings of the same fact: $\tfrac12$ (modulus, spin,
fixed point of $1-x$); the $0|1$ of the fermionic superpoint; and the tension $T_{AB}$ that
produced it.

From this object follows its one-dimensional description: the worldline and its path
integral, the Gaussian normalised by the half-factorial,
$\int_{-\infty}^{\infty}e^{-x^2}dx=\sqrt\pi$, $(\tfrac12)!=\Gamma(\tfrac32)=\tfrac{\sqrt\pi}{2}$
(\ProofWitness{gamma\_half}), computed where it is stated. The particle is then a type~I
operator: real discrete spectrum, minimal projections, integer dimension. Its internal
structure is the spin-star,
\begin{equation}\label{eq:obs-spinstar}
  \underbrace{S+E_1+\cdots+E_N}_{\text{central}+\text{environment}}
   \xrightarrow{\;N=2\;}\underbrace{C+P+F}_{\text{PCF}},\qquad \Fmax=N^{2}=4 .
\end{equation}
with arithmetic $1+2=3$: the binary fermion, mediated by $\tfrac12$, has become the
tripartite particle. And the geometry that generated the particle fixes its charges: the
eigenvalue triangle is the Eisenstein lattice $\Z[\omega]=A_2=\mathfrak{su}(3)$; the Clifford
latitude gives $SU(2)$; the Hopf fibre, $U(1)$; the weak mixing angle at
unification is a ratio of mode counts, $\sin^2\theta_W|_{\rm GUT}=N(0)/N(2)=\tfrac38$, which runs
to the measured $0.231$ at $M_Z$ through the MSSM $\beta$-functions
(\S\ref{app:gauge}; \ProofWitnessBlock{gauge\_dim\_su3, hopf\_from\_clifford, weinberg\_angle\_gut}).
\begin{figure}[!htbp]\centering
\includegraphics[height=0.82\textheight,keepaspectratio]{fig6_cylinder_torus.pdf}
\caption{From $\golden^{2}=\golden+1$: (a) the cylinder $C_0$ with $\mathrm{P,C,F}$ at $120^\circ$
($\|P\|\,\|C\|\,\|F\|=\tfrac12$) closes to (b) the torus $T^2_{\PCF}=\C/\Lambda$, $\tau=i$
($|z_1|^2+|z_2|^2=\tfrac14{+}\tfrac34=1$, the $S^3$ constraint), whose
Clifford embedding $T^2\hookrightarrow S^3$ and Hopf fibre give the gauge factors
$\mathrm{SU}(3)\times\mathrm{SU}(2)\times\mathrm{U}(1)$ of App.~\ref{app:gauge}.}
\label{fig:torus-gauge}
\end{figure}
\paragraph{Geometry and particle are the same object.}
The order of these facts is fixed by the diagram. The
projection $\Pi\colon\E^3\to\C$ is the \emph{root} of the whole downstream chain, and it is
at once the geometry of spacetime---the sheet and torus projected---and the particle. This
reflects the \emph{self-similarity} of the object: the same modulus $\tfrac12$ at every scale
level, seeded by $\golden^2=\golden+1$ extending $i^2=-1$ from the plane to space. The
particle carries the same geometry in its information and, as it accumulates entropy, keeps
it and moves within the spacetime geometry it itself is. There is no stage with geometry and
a later one with the particle: they are the same object at two scale readings. It is this
accumulation of information, carried by the self-similar geometry, that the next movement
reads as an observer (\S\ref{ssec:observer}): container and content, bulk and boundary,
geometry and particle are not separate---the same non-separation the M\"obius band already
fixed.

\emph{W1 --- the observer described by the theory (CLPW).} CLPW require (\S2.5) that the
observer be described by the theory, not introduced from outside. $\Pi$ is internal---no
external region from which to look---and self-similar, so the observer is the geometry itself
read at its centre. \emph{Proof:} \eqref{eq:obs-interface} with the bulk--boundary isometry
of \S\ref{sec:derivations}; Proposition~\ref{prop:obs-nofirewall};
\cite{CLPW22}. \ProofWitnessBlock{cw3\_observer\_items\_core}

\emph{W2 --- finite-dimensional Hilbert space (W01).} The finite horizon area requires finite
dimension; here the spectrum is the triangle of three eigenvalues and the mode count is
integer. \eqref{eq:obs-spinstar}. \ProofWitnessBlock{witten\_finite\_rank}

\subsection{The particle as observer}\label{ssec:observer}

What turns the particle into an observer is the accumulation of information, described
here by an internal Wick rotation, not imposed. When the centre is taken
as the time axis, its square changes sign, $(i\,t)^2=-t^2$ (\ProofWitness{wick\_squares\_flip\_sign, lorentzian\_signature\_sum}), and past and future become the two light cones, with the
apex---the present---as observer (\ProofWitness{C\_fixed\_PF\_rotate, causal\_trichotomy}). The
full metric is formalised in \S\ref{ssec:accum}; here the algebra of the rotation, which
is self-contained, suffices. The observer accumulates: its Fisher information grows to the
ceiling the Cram\'er--Rao bound permits,
\begin{equation}\label{eq:obs-fishertime}
  \tau_F=\frac{\tau_D}{\sqrt{2f}},\qquad \tau_F=\tau_D\iff f=\tfrac12 .
\end{equation}\begin{equation}\label{eq:obs-cramerrao}
  \mathrm{Var}(\theta)\ge F_\theta^{-1},\qquad
  \underbrace{\Fmax^{-1}=\tfrac14}_{\text{Kiely: Cram\'er--Rao}}
   \;=\;\underbrace{|\Om|^{2}=\mu_3^{2}=\tfrac14}_{\text{PCF area factor}} .
\end{equation}\begin{equation}\label{eq:obs-redundancy}
  R_\delta\sim N\ \ (\text{all fragments agree})\;\longleftrightarrow\;\text{objective record},
\end{equation}\begin{equation}\label{eq:obs-accum}
  F_\theta(t)\xrightarrow{\;N\to\infty\;}\Fmax\big(1-e^{-(t/\tau_F)^{2}}\big),\quad \Fmax=4,
  \qquad F_\Om=4\mu_3^{2}=1\ \text{bit},
\end{equation}
and sees exactly the half,
\begin{equation}\label{eq:obs-half}
  |P|\cdot|C|\cdot|F|=\tfrac12=|\Om| .
\end{equation}
The ``half'' is geometric: the eigenvalue dipole places $P$ at the south pole and $F$ at
the north, with $C$ the equatorial axis ($|C|=1$), and $|P|\,|F|=\tfrac14$ while $|C|=1$
gives the product $\tfrac12$; the $\sqrt3$ of $|F|$ cancels that of $|P|$---a cancellation
proper to the equilateral triangle, arity three (Proposition~\ref{prop:pcf-norms},
Proposition~\ref{prop:spectral}: $\sigma+\mu=2$, $\sigma\mu=\tfrac34$ with unique solution
$\mu=\tfrac12$, $\sigma=\tfrac32$.

\emph{W5 --- the initial--final pairing observable from within (W01).} Witten computes the
pairing $\langle f|i\rangle$ as calculable but not observable from outside (W01). It
becomes the internal pairing $\langle F|P\rangle=\bar F P/C$, with
$|\langle F|P\rangle|=\tfrac12$, observable because $\Pi$ is inside and no external region
remains. \emph{Proof:}

\emph{W3 --- the unitarity conjecture (W01).} Witten conjectures that the Hermitian
pairing is positive definite (W01). The spectrum is positive and $|\Om|=\tfrac12>\tfrac14$, so the object does not collapse.
\emph{Proof:}
(4.7).

The coincidence with Kiely's quantum Darwinism is forced, not numerical. The objectivity
threshold is $f_{\rm crit}=\tfrac12$; the redundancy grows as the mode number; the
certainty condition is $\varepsilon_0 M=\pi$:
\begin{equation}\label{eq:obs-threshold}
  \underbrace{f_{\mathrm{crit}}=\tfrac12}_{\text{Kiely: objectivity}}
   \;=\;\underbrace{|\Om|=\mu_3=\tfrac12}_{\text{PCF}},
\end{equation}\begin{equation}\label{eq:obs-certainty}
  \epsz\cdot\Mpcf=\pi\qquad(\text{cell capacity }=\pi\text{ bits}),
\end{equation}
\begin{proposition}[No separable parts before observation]\label{prop:no-parts}
For a pure two-party state let $p_1\ge p_2\ge\dots\ge0$ be the eigenvalues of the reduced state,
$\sum_i p_i=1$ (these are the squared Schmidt coefficients, $p_i=\lambda_i^2$; the criterion below
holds in this convention and not for the $\lambda_i$ themselves). The entanglement modulus is the
top eigenvalue, $|\Om|=p_1$, and it satisfies $|\Om|=|\Om|^{2}$ exactly when $p_1=1$---that is,
when every other $p_i$ vanishes, the reduced state is pure, the Schmidt rank is one and the state is
a product---and $|\Om|>|\Om|^{2}$ otherwise. The No-Diagonal fixes
$|\widehat\Om|=\tfrac12$ in the balanced (maximally entangled) case
(Proposition~\ref{prop:spectral}), so
\[
  |\Om|=\tfrac12\;>\;\tfrac14=|\Om|^{2},\qquad H(\tfrac12)=1\ \text{bit},
\]
the strict inequality being the signature of non-factorizability; for a product state it collapses to
$|\Om|=|\Om|^{2}=1$, $H=0$  (the Schmidt decomposition~\cite{Schmidt1907}) . Hence before observation there are no separable subsystems carrying \ProofWitness{no\_separable\_parts, separable\_iff\_top\_eigenvalue\_one}
values: what is real is the non-local whole, recovered as $|\widehat\Om|=\tfrac12$. PCF therefore
predetermines no parts, so it is neither a local-realist (Bell) nor a parts-realist (Leggett)
hidden-variable model, and the no-go theorems of~\cite{Bell64} and~\cite{Leggett03} do not bind it. 
\end{proposition}

The No-Diagonal theorem (\S\ref{ssec:arity}) establishes that every self-referential
structure without collapse must produce $\mu^2=\tfrac14$ and the tripartite
decomposition, with no other solution---which forces each identification. In the diagram,
here the first apex converges: three independent routes---the bridge (\S\ref{ssec:hinge}),
the triangle norm and the Fisher metric---give the same $\tfrac12$. When three independent
chains reach the same number, the diagram commutes at that point.

\subsection{The tower of information, and the emergence of curved spacetime}\label{ssec:accum}

\paragraph{The tower (types I and II interwoven).}
The accumulating observer does not select one reality among many: it builds it. There is
no landscape of vacua to choose from, no statistical ensemble to average. This
homogeneity---the same structure at every point---is what the field equations require: it
is the tensor covariance of \S\ref{ssec:hinge} now read as a constructive principle, and
its full development comes below. The entropy of one bit per cell, $H(\tfrac12)=1$
(\ProofWitness{binary\_entropy\_half}), is the maximally mixed state---$\rho_{\max}=1$ in CLPW's
language---describing the de~Sitter vacuum, and from its overflow the geometry is born.

When the $\pi$-bit cell fills, the throat opens:
\begin{equation}\label{eq:obs-throat}
  S(\sigma)=\pi\golden^{\sigma},
\end{equation}
with $N_{\rm modes}(\sigma)=\lfloor\pi\golden^\sigma\rfloor$. As it overflows, the throat
count is set:
\begin{equation}\label{eq:obs-swampland}
  \frac{|\partial_\sigma V|}{V}=\ln\golden=c_{\mathrm{PCF}} .
\end{equation}
The UV fixed point forces the invariants:
\begin{assumption}[The gauge--tower bridge]\label{asm:bridge}
\begin{equation}\label{eq:obs-fixedpoint}
  \beta_g(g^{*},\lambda^{*})=0\;\Longleftrightarrow\;\epsz\Mpcf=\pi .
\end{equation}
The right-hand side is derived and triply proved---in \S\ref*{ssec:construction} as
\eqref{eq:certainty}, here as \ProofWitness{eps0\_Mpcf\_pi}, and checked numerically \NumericalWitness{eq:certainty}; the equivalence itself is
this section's one input, the bridge between the gauge running and the tower's algebra. \Conj
\end{assumption}
The quantisation is in powers of $\golden$, and the identity that closes it is
$\golden^{-\lambda}=\tfrac12\Leftrightarrow\golden^\lambda=2$: the binary base comes out of
the golden geometry. Here the string reappears, not as a new object but as the same
particle in its information register: the dimension of the state space at each level is
$\dim\mathcal{H}(\sigma)=\lfloor S(\sigma)\rfloor$.

\emph{W4 --- the entropy conjecture (W01).} Witten conjectures that the Hilbert space is
finite-dimensional with $S=\ln N$ (W01). The Hilbert dimension per level is the integer
part of the horizon entropy, $\dim\mathcal{H}(\sigma)=\lfloor S(\sigma)\rfloor$, with
$N_{\rm modes}=\lfloor\pi\golden^\sigma\rfloor$---the constructive form of Witten's
$S=\ln N$. \emph{Proof:}
\ProofWitnessBlock{witten\_entropy\_realized, witten\_finite\_rank, witten\_two\_conjectures} (by \texttt{rfl}---the most direct realisation).
The entropy
branch is asserted strongly (tight); the unitarity branch is ``realised''.

The flow carrying one level to another, $\sigma\mapsto\sigma+t$, fixes $\tau=i$ and
preserves $\tfrac12$: it is the tracial flow, the type~II signature, and it is the thermal
flow of the static patch at the Gibbons--Hawking temperature.
\ProofWitnessBlock{tower\_scale\_group, tower\_flow\_fixes\_tau, omega\_flow\_group, omega\_flow\_invariant}

\emph{W11 --- the type~II algebra by crossed product (CLPW, completed).} CLPW build the
type~II$_1$ algebra from the type~III by the crossed product with the modular flow
(\S1.2--1.3). The scale flow preserves $\tfrac12$ as a trace, which is the type~II
signature; the section proves the flow and takes from CLPW the name of the crossed-product
classification. \emph{Proof:}
\ProofWitnessBlock{omega\_flow\_invariant, tower\_scale\_group}.

\emph{W12 --- the maximally mixed vacuum state (CLPW \S2.4, developed).} CLPW identify the
de~Sitter vacuum with the maximally mixed state, $\rho_{\max}=1$. One bit per cell,
$H(\tfrac12)=1$, gives that $\rho=1$ of maximal entropy. \emph{Proof:}
\ProofWitnessBlock{binary\_entropy\_half}.

The heart of the movement is that bulk and boundary are the same thing at once. The bulk
is the three-dimensional projection; the boundary, the two-dimensional plane; and they are
the same $\Pi$ of the first movement. The two accumulations---of time (Fisher) and of
scale (tower)---are one:
\begin{equation}\label{eq:obs-weld}
  \tau_F(\sigma)=\tau(\sigma)=\Mpcf\,\golden^{-\sigma},\qquad
  S(\sigma)\,\tau_F(\sigma)=\pi\Mpcf\ \ (\text{constant in }\sigma).
\end{equation}
The two lengths the tower carries are conjugate: the throat coordinate $z(\sigma)=\golden^{\sigma}$
grows with the level and the weld time $\tau(\sigma)=\Mpcf\golden^{-\sigma}$ shrinks with it, and
their product is constant,
\begin{equation}\label{eq:conjugate-pair}
  z(\sigma)\,\tau(\sigma)=\Mpcf\qquad\text{for every }\sigma .
\end{equation}
This is the relation $R\cdot(\alpha'/R)=\alpha'$ with $\alpha'=\Mpcf$: the compactification radius is
the throat coordinate, the lattice spacing is its T-dual, and $\alpha'$ is \emph{fixed by the product}
rather than chosen. Short and long distance are therefore the same family exchanged by duality, not two
regimes with a limit between them---the structural reason there is no ultraviolet regulator to remove
(Remark~\ref{rmk:residual}). \ProofWitnessBlock{lattice\_radius\_conjugate}

\begin{equation}\label{eq:obs-identity}
  \underbrace{F_\Om\cdot\Nmodes}_{\text{bits (Shannon)}}
   =\underbrace{S(\sigma)=\pi\golden^{\sigma}}_{\text{tower}}
   =\underbrace{\Nmodes=\lfloor S(\sigma)\rfloor}_{\text{matter}}
   =\underbrace{\SBH=\frac{A}{4\GN}}_{\text{horizon}}
   =\underbrace{S_{\text{Jacobson}}}_{\delta Q=T\delta S},
\end{equation}
This identification---the \emph{weld}---is where four routes converge, and it is the only
point where the three types coincide in one object: type~I by the discrete tower, type~II
by the observer's time, type~III because it lives on the spacetime. It is the maximal apex
of the diagram. Local and global are not separate: $\tfrac12$ at each level (local) and
invariant under the flow (global), the same invariant read two ways. Its last two branches
close below.

\paragraph{Curved spacetime (type III again).}
The accumulated entropy now sources gravity. Each bit costs---the Landauer cost, which is the
regulator $R_K=\log\golden$ of $\Q(\sqrt5)$ (Definition~\ref{def:regulator}) measured in
$\lambda_{\log}$, Theorem~\ref{thm:entropy-bridge},
$\lambda\cdot R=\log 2$---and when that cost is imposed as an equation of state
$\delta Q=T\,\delta S$ across every local Rindler horizon, the result is, by Jacobson's
argument, the Einstein equation:
\begin{equation}\label{eq:obs-landauer}
  \text{energy per bit}=\frac{1}{\Mpcf},\qquad
  \SBH/k_B=\lng\cdot R_K=\log2\ \ (\text{one bit; Theorem~\ref{thm:entropy-bridge}}),
\end{equation}\begin{equation}\label{eq:obs-jacobson}
  \delta Q=T\,\delta S\ \Longrightarrow\ \text{the Einstein equation},
\end{equation}\begin{equation}\label{eq:obs-einstein}
  R_{AB}=-4\,g_{AB},\quad R=-20,\qquad G_{AB}+\Lambda_5\,g_{AB}=\kappa\,T_{AB},\quad
  \kappa=8\pi\GN,
\end{equation}
with $G_N=\mu_3=\tfrac12$. Here the homogeneity of \S\ref{ssec:hinge} is redeemed:
Jacobson's argument requires a local horizon with the same law at every point---exactly
what tensor covariance guarantees; without it the step does not exist. The matter is the
mode number:
\begin{equation}\label{eq:obs-matter}
  T^{\mathrm{YM}}_{AB}=F_{AC}F_B{}^{C}-\tfrac14\,g_{AB}F_{CD}F^{CD},\qquad
  \text{matter}=\Nmodes=\lfloor S(\sigma)\rfloor .
\end{equation}
and here the two open branches of the identity close: the horizon entropy and Jacobson's
were, all along, the same number as the bits and the modes. Each mode is one bit exactly at the
modulus: the per-eigenvalue ceiling $F_\Om=4\mu_3^{2}$ is unity iff $\mu_3=\tfrac12$, and the
binary entropy there is $H(\tfrac12)=1$ bit.
\ProofWitnessBlock{fisher\_unit\_iff, H\_half\_one\_bit}%

That matter does not merely sit on the background: it curves it. The back-reaction is fixed by the
same modulus, with no free parameter, and it is what makes the shared root of \emph{metric} and
\emph{tension} (\S\ref{ssec:hinge}) a computation rather than an etymology.

\begin{definition}[Shell tension]\label{def:shell-tension}
A thin matter shell at the discrete level $w_\sigma=\sigma\log\golden$ of the throat
\eqref{eq:throat} has \emph{tension} $\lambda_\sigma$: its energy density per unit volume of the
constant-$w$ leaf. Its value is pinned by Proposition~\ref{prop:shell-tension}.
\ProofWitnessBlock{prop:shell-tension}%
\end{definition}

\begin{proposition}[The energy per bit is constant along the tower]\label{prop:ebit}
For every level,
\begin{equation}\label{eq:ebit}
  \frac{\varepsilon(\sigma)}{S(\sigma)}=\frac{\epsz}{\pi}=\frac{1}{\Mpcf},
\end{equation}
equivalently $S(\sigma)=\Mpcf\,\varepsilon(\sigma)$: the first law with temperature $T=1/\Mpcf$.
\ProofWitnessBlock{energy\_per\_bit\_constant}
\end{proposition}
\begin{proof}
The per-level energy is $\varepsilon(\sigma)=\epsz\golden^{\sigma}$ \eqref{eq:obs-accum} and the
saturation entropy is $S(\sigma)=\pi\golden^{\sigma}$ \eqref{eq:tower-modes}; the quotient cancels
$\golden^{\sigma}$ and leaves $\epsz/\pi$, which is $1/\Mpcf$ by the certainty relation
\eqref{eq:certainty}. The cancellation is what the common base buys: an energy growing at any other
rate would leave a $\sigma$-dependent residue.
\end{proof}

\begin{proposition}[The shell tension carries no free parameter]\label{prop:shell-tension}
The tension of Definition~\ref{def:shell-tension} is fixed by the level's mode count,
\begin{equation}\label{eq:shell-tension}
  \lambda_\sigma=\frac{\Nmodes(\sigma)}{\Mpcf}.
\end{equation}
\ProofWitnessBlock{shell\_tension\_determined}
\end{proposition}
\begin{proof}
Each tower mode is one bit: $\Nmodes(\sigma)=\lfloor S(\sigma)\rfloor$ \eqref{eq:obs-matter} is one
mode per bit of horizon entropy, and the Kiely ceiling $F_\Om=4\muthree^{2}=1$
\eqref{eq:obs-identity} makes each mode exactly one bit, neither more nor less. The shell's tension
is then the energy per bit times the number of bits,
$\lambda_\sigma=\varepsilon(\sigma)\Nmodes(\sigma)/S(\sigma)$, which is
$\Nmodes(\sigma)/\Mpcf$ by Proposition~\ref{prop:ebit}.
\end{proof}

\begin{proposition}[Israel junction: the back-reaction of each level]\label{prop:israel}
A constant-$w$ leaf of \eqref{eq:throat} has extrinsic curvature $K_{\mu\nu}=A'(w)g_{\mu\nu}$, so a
shell of tension $\lambda_\sigma$ imposes the jump
\begin{equation}\label{eq:israel}
  [A'(w_\sigma)]=-\frac{8\pi G_5}{3}\,\lambda_\sigma=-\frac{4\pi}{3}\,\lambda_\sigma,
  \qquad G_5=\muthree,
\end{equation}
and with \eqref{eq:shell-tension} the jump is determined level by level,
$[A'(w_\sigma)]=-\tfrac{4\pi}{3}\Nmodes(\sigma)/\Mpcf$. Between shells the metric is exact
$\mathrm{AdS}_5$: a steeper warp is a deeper curvature.
for the arithmetic; the junction condition itself~\cite{Israel66}. \ProofWitnessBlock{israel\_prefactor, israel\_jump\_determined}

\end{proposition}
\begin{proof}
The five-dimensional junction condition for a thin shell is
$[K_{\mu\nu}]=-8\pi G_5\bigl(S_{\mu\nu}-\tfrac13h_{\mu\nu}S\bigr)$; for pure tension
$S_{\mu\nu}=-\lambda h_{\mu\nu}$ this reduces to $[A']=-\tfrac{8\pi G_5}{3}\lambda$. With
$G_5=\muthree$ \eqref{eq:brown-henneaux} the prefactor collapses to $4\pi/3$---and only there: any
other Newton constant leaves $8\pi G_5/3\neq4\pi/3$.
\end{proof}

\begin{remark}[The cumulative back-reaction]\label{rmk:backreaction}
Each level kinks the warp---and deepens the local $|\Lambda|$---by an amount set by its mode count,
so the cumulative back-reaction up to level $k$ is $\sum_{\sigma\le k}\Nmodes(\sigma)$, giving
$3,8,16,29,50,84,140$ for $k=0,\dots,6$. The gravitational side and the string side are here the
same computation: \eqref{eq:israel} ties the extrinsic curvature to the shell tension, and
\eqref{eq:shell-tension} ties that tension to the mode count, so what curves the geometry is the
number of bits the level carries.
\ProofWitnessBlock{backreaction\_cumulative}
\NumericalWitnessBlock{eq:shell-tension}
\end{remark}

\noindent The cosmological constant is determined, not chosen:

\emph{W8 --- a theory that determines $\Lambda$ (W01).} The abstract of W01 states that
general relativity must be derived from a theory that \emph{determines} the value of
$\Lambda$. Here $\Lambda$ comes from the curvature, with zero free parameters:
$\Lambda_5=-d(d-1)/2\ell^2=-6$, with $\ell=1$ from \eqref{eq:bridge-S} (\S\ref{sec:derivations})
and $d=4$ from the arity (\S\ref{sec:methods}), and the G--$\Lambda$ duality
$\golden^{-6}\cdot\golden^{6}=1$ (\ProofWitness{G\_Lambda\_duality, R\_scalar\_pcf, BF\_value}).
Between the gravity scale and the $\Lambda$ scale there is an exact interval:

It remains to show the spacetime is de~Sitter. The metric obtained is Euclidean
anti-de~Sitter; the same internal Wick rotation that made the centre the time axis carries
it to de~Sitter, with $\varepsilon(\sigma)=\varepsilon_0\golden^\sigma$ and $\lambda$ the length
that converts the tower's logarithmic step into proper length along the scale axis, so that
$X_4=\lambda\ln s$ is the fifth, flat coordinate of \S\ref{app:embedding}:
\begin{equation}\label{eq:ets-metric}
  ds^{2}=\underbrace{(dx^{2}+dy^{2}+dz^{2})}_{\E^{3}}
        -\underbrace{c^{2}\,dt^{2}}_{\text{Wick }(C)}
        +\underbrace{\lambda^{2}\,d(\ln s)^{2}}_{\golden^{\sigma}\text{-tower}},
  \qquad s=\golden^{\sigma},
\end{equation}
It is the bridge to de~Sitter that Maldacena anticipated, built here by internal analytic
continuation rather than an imposed duality (\ProofWitness{ets\_is\_ds\_zero\_h}). Past and future
are the poles of a dipole, with $C$ at the equator---the two conjugate sheets of
de~Sitter.

\emph{W10 --- the antipodal map / meta-observables (W01).} The antipodal singularity
Witten notes corresponds to $F$ being antipodal to $P$ in the dipole, which resolves
Strominger's obstruction geometrically. \emph{Proof:} the dipole of \eqref{eq:ets-metric};
\cite{Witten01dS} (\S\ref{app:embedding}).

The geometric closure is the connection the framework contributes. The 5D ambient is
flat---its curvature, computed, vanishes, $R^{(5)}=0$ (\ProofWitness{ets\_constant\_block\_christoffels\_vanish})---and the
fifth coordinate, the scale axis, is the extra direction an embedding needs. De~Sitter is
the hyperboloid $-X_0^2+\cdots+X_4^2=1/H^2$ embedded in that flat ambient, with intrinsic
curvature ${}={}$ wedge of the extrinsic by Gauss: $R_{\mu\nu}=3H^2 g_{\mu\nu}$
(\ProofWitness{ds\_ricci\_from\_gauss, ds\_einstein\_lambda}). The planar patch covers exactly half
the hyperboloid---$X_0+X_4=\ell\,e^{t/\ell}>0$ never changes sign
(\ProofWitness{ds\_covers\_half\_hyperboloid})---and that half \emph{realises} the observer's half
$|\Om|=\tfrac12$ by the shared value (\ProofWitness{observer\_half\_from\_norms}). The assertoric
level splits: strong on the covering, ``realises'' on the identification.

\emph{W9 --- the thermal nature of the static patch (Gibbons--Hawking).} This is the
classical Gibbons--Hawking result, prior to W01 and CLPW: the de~Sitter horizon radiates
and the static patch is a thermal state at $T=H/2\pi$. The observer's modular flow is the
tracial flow of the tower, and the static patch is thermal at that temperature.
 (\S\ref{app:embedding}) \cite{GibbonsHawking77}. \ProofWitnessBlock{dS\_temperature, ds\_beta\_reciprocal}

\medskip
The formal statements underlying the accumulation and gravity movements follow; their
proofs are those of the corresponding \ProofWitness{...} tags.
\begin{remark}[The role of $\golden$ in the observer]
$\mu_3=\tfrac12=\golden^{-\lng}$; the three invariants ($f_{\mathrm{crit}}=\tfrac12$,
$\Fmax^{-1}=\tfrac14$, $n=3$) are $\golden$-facts of \S\ref{sec:methods}--\S\ref{sec:derivations},
not fitted values.
\end{remark}\begin{remark}[The role of $\golden$ in the objectivity threshold]
$|\Om|=\tfrac12=\mu_3=\golden^{-\lng}$: the half of de Sitter, Kiely's threshold and the modulus
are one $\golden$-fact; the throat that observation builds is the tower $\golden^{\sigma}$.
\end{remark}\begin{remark}[The role of $\golden$ in the accumulation]
$\epsz=\ln\golden/(6\sqrt3)$ and $\Mpcf=6\sqrt3\,\pi/\ln\golden$ give $\epsz\Mpcf=\pi$: the clock
of the accumulation, the ultraviolet fixed point and the certainty principle are one
$\golden$-fact; the accumulation tower is $\golden^{\sigma}$.
\end{remark}\begin{remark}[Past and future are fixed dynamically]
The reading of $C$-$P$-$F$ as present\,/\,past\,/\,future is not a free labelling but a consequence
of the causal structure. The Wick rotation makes the centre $C$ the time axis and gives the metric
of \eqref{eq:ets-metric} its Lorentzian signature $(+,+,+,-)$, so it carries a light cone: $C$ is
its apex---the observer's present---and $P,F$ are the past- and future-directed cones, rotating in
opposite senses about the fixed, real $C$, with $\Theta$ exchanging them. The direction of the
$\golden$-driven accumulation along the throat---from the origin toward the ceiling
(Proposition~\ref{prop:obs-accum})---orients the cone, singling out one direction as the future.
The past/future assignment therefore follows from the causal cone and the direction of
accumulation, not from convention.
\end{remark}\begin{proposition}[Accumulation and ceiling]\label{prop:obs-accum}
The metrological quantum Fisher information accumulates as
$F_\theta(t)\to\Fmax\bigl(1-e^{-(t/\tauF)^{2}}\bigr)$ with $\Fmax=4$~\cite{Kiely26},
so the Cram\'er--Rao ceiling $\Fmax^{-1}=\tfrac14=\muthree^{2}$ equals the
holographic area factor. The Fisher information per eigenvalue is
$F_{\Om}=4\muthree^{2}=1$ bit; the three eigenvalues give $3F_{\Om}=3=c$, the
central charge of \S\ref{ssec:object}; and the objectivity threshold
$f_{\mathrm{crit}}=\muthree=\tfrac12$ is the maximum-entropy point $H(\tfrac12)=1$ bit
(Proposition~\ref{prop:entropy-max}).
\ProofWitnessBlock{binary\_entropy\_half, S\_tower\_recurrence}
\NumericalWitnessBlock{prop:entropy-max}
\end{proposition}

\begin{theorem}[Time--scale conjugacy]\label{thm:obs-weld}
At the operating point $f=\muthree=\tfrac12$ the metrological timescale equals the
tower period, $\tauF(\sigma)=\tau(\sigma)=\Mpcf\,\golden^{-\sigma}$. Hence
\[
S(\sigma)\,\tauF(\sigma)=(\pi\golden^{\sigma})\,(\Mpcf\,\golden^{-\sigma})=\pi\Mpcf
\]
is constant in $\sigma$: the certainty principle $\epsz\Mpcf=\pi$
(\eqref{eq:certainty}) promoted from a single cell to the generation/clock pair.
\ProofWitnessBlock{obs\_weld}
\end{theorem}\begin{proposition}[The bulk metric is an Einstein space]\label{prop:obs-einstein}
For the warped throat \eqref{eq:throat} with $A=-\log z$ (anti-de~Sitter,
$\ell=1$), explicit computation of the Levi-Civita curvature gives, uniformly in
all five components,
\[
R_{AB}=-4\,g_{AB},\qquad R=-20,\qquad G_{AB}=6\,g_{AB}\quad(A,B=0,\dots,4),
\]
so the metric solves $G_{AB}+\Lambda_5\,g_{AB}=0$ with $\Lambda_5=-6$, and the
Breitenlohner--Freedman bound $m^{2}_{\mathrm{BF}}=-d^{2}/4=-4$ holds, so the discrete tower is
stable. With a
source, $G_{AB}+\Lambda_5\,g_{AB}=\kappa\,T_{AB}$ and $\kappa=8\pi\GN$; the PCF
value $\GN=\muthree=\tfrac12$ gives $\kappa=4\pi$ (effective $G_4=\tfrac14$), and
$\nabla^{A}G_{AB}=0$ holds identically, so $T_{AB}$ is conserved.
\ProofWitnessBlock{G\_AB\_pcf, R\_AB\_einstein, R\_scalar\_pcf, einstein\_coeff, ricci\_coeff, ricci\_scalar, BF\_value}
\end{proposition}

\begin{theorem}[Jacobson's construction on the bulk metric]\label{thm:obs-jacobson}
Since \eqref{eq:throat} is an Einstein space, $R_{ab}k^{a}k^{b}=-4\,g_{ab}k^{a}k^{b}=0$
for every null $k$; in vacuum the Clausius relation
$\delta Q=\int T_{ab}\chi^{a}\,d\Sigma^{b}=0$ forces $\delta S=0$ (Jacobson's
entanglement equilibrium~\cite{Jacobson1995}), which coincides with the ceiling
$\Fmax^{-1}=\muthree^{2}$ of Prop.~\ref{prop:obs-accum}. The normalisation
$\eta=1/4\GN$ with $\GN=\muthree$ fixes the area factor $\muthree^{2}=\tfrac14$
and the coupling $2\pi/\eta=8\pi\GN$; with matter, $\delta Q=T\,\delta S$ on every
local Rindler horizon yields the Einstein equation~\cite{Jacobson1995,Jacobson2016}.
\ProofWitnessBlock{ds\_einstein\_lambda, R\_AB\_einstein, G\_AB\_pcf}
\end{theorem}

\begin{proposition}[Matter is the saturation entropy]\label{prop:obs-matter}
Each interior level carries $\Nmodes(\sigma)=\lfloor\pi\golden^{\sigma}\rfloor
=\lfloor S(\sigma)\rfloor$ scalar modes, each contributing one Shannon bit, so the
matter content of level $\sigma$ is the saturation entropy $S(\sigma)=\pi\golden^{\sigma}$. That
this is an \emph{equality} and not merely a count is the first law of the tower: the Kiely ceiling
$F_\Om=4\muthree^{2}=1$ makes each mode exactly one bit (Proposition~\ref{prop:obs-accum}), the
per-level energy is $\varepsilon(\sigma)=\epsz\golden^{\sigma}$ with $\epsz\Mpcf=\pi$, so the energy
per bit is constant across the tower (Proposition~\ref{prop:ebit}),
\begin{equation}
  \frac{\varepsilon(\sigma)}{S(\sigma)}=\frac{\epsz}{\pi}=\frac1{\Mpcf},
\end{equation}
which is $S(\sigma)=\Mpcf\,\varepsilon(\sigma)$ at temperature $T=1/\Mpcf$
The colour octet has $\dim\mathfrak{su}(3)=n^{2}-1=8$ at the arity
$n=3$ fixed by the golden central chain $\golden^{2}+\golden^{-2}=3$; the full
$\mathrm{SU}(3)\times\mathrm{SU}(2)\times\mathrm{U}(1)$ follows from the torus $T^2_{\PCF}$
\eqref{eq:torus}, its gauge factors landing at the force levels by increasing dimension
$1<3<8$: $\mathrm{U}(1)$ is the Hopf fibre $S^1$ at $\sigma=3$ (EM), $\mathrm{SU}(2)$ the Clifford
$S^3$ at $\sigma=4$ (weak), and $\mathrm{SU}(3)$ the $A_2$ roots of the Eisenstein face
(Fig.~\ref{fig:isometry-algebra}) at $\sigma=5$ (strong), with $\dim\mathfrak{su}(3)=n^{2}-1=8$ at
the arity $n=3$ fixed by the golden central chain $\golden^{2}+\golden^{-2}=3$
\ProofWitnessBlock{gauge\_dim\_su3, phi\_central\_chain, central\_chain, KK\_numerator, clifford\_S3\_condition, hopf\_from\_clifford, a2\_minimal\_vectors\_exactly\_six}.

The electroweak data follow from the same modulus: the weak mixing angle at unification is
$\sin^{2}\theta_W|_{\rm GUT}=\tfrac38=N(0)/N(2)$ \ProofWitness{weinberg\_angle\_gut}
running to $0.231$ at
$M_Z$; the tower entropy ratio $S(3)/S(6)=\golden^{-3}$ \ProofWitness{entropy\_ratio\_S3\_S6} marks the
breaking $\Z_3\to\Z_4$ at the tower midpoint $\sigma=3$, and the $G$--$\Lambda$ duality reads $\golden^{-6}\golden^{+6}=1$
The generations are not placed radially by this $\sigma$-assignment:
the charged-lepton ratios are not integer powers of $\pi$ ($m_\mu/m_e=\pi^{4.66}$), and a
$\pi$-graded radial triple cannot satisfy the Koide relation. Two statements are at work here, of different objects: the Koide ratio $Q=\tfrac23$ is an algebraic identity of that placement form, holding for every $\delta$ and hence no evidence for any particular value; the prediction is the three masses themselves, reconstructed to $\sim10^{-4}$ from the single scale $\sqrt M$ with $\delta=2/n^{2}$ fixed by the arity, a value so sharply selected that a $0.1\%$ departure degrades the fit by two orders of magnitude. The $\sigma$-levels label sectors and
forces; the three generations are the \emph{angular} $\omega$-triple within a sector, placed in the
bridge with no free angular parameter,
$\sqrt{m_k}=\sqrt M\,(1+\sqrt2\cos(\delta+2\pi k/3))$ with $\delta=2/n^{2}=\tfrac29$ at arity $n=3$,
reconstructing the three charged-lepton masses to ${\sim}10^{-4}$ from the single scale $\sqrt M$.
That scale is not an independent input: the certainty condition $\varepsilon_0 M=\pi$ \eqref{eq:obs-threshold} is the same that fixes $\Mpcf$, so $M=\Mpcf$; Theorem~\ref{thm:colour-from-M} then derives the colour scale and running from it. What remains open is only the quark spectrum and its CKM mixing, not the scale or the colour running.   (bridge)% \ProofWitness{M\_eq\_Mpcf}
Their Yang--Mills stress-energy
$T^{\mathrm{YM}}_{AB}=F_{AC}F_{B}{}^{C}-\tfrac14\,g_{AB}F_{CD}F^{CD}$ is
conserved ($\nabla^{A}T^{\mathrm{YM}}_{AB}=0$) and traceless on the
$D=4$ worldsheet (conformal there).
\end{proposition}

\begin{theorem}[The Landau--Lifshitz energy of the field]\label{thm:LL-energy}
On the de~Sitter background \eqref{eq:ets-metric} with $\GN=\muthree=\tfrac12$, the Landau--Lifshitz complex $(-g)(T^{\mu\nu}+t_{\mathrm{LL}}^{\mu\nu})=\tfrac1{16\pi\GN}\partial_\alpha\partial_\beta H^{\mu\alpha\nu\beta}$ has vanishing energy density---its $00$ component is zero, the net energy of matter and field balancing, the classical face of the entanglement equilibrium $\delta S=0$ of Theorem~\ref{thm:obs-jacobson}---while its spatial components $-2H^{2}e^{4Ht}/\pi$ do not vanish, as de~Sitter is not static (the standard complex differs from the physical vacuum pressure by the $\Lambda$-term the pseudotensor requires for $\Lambda\neq0$). The gravitational energy is a definite boundary charge: the energy in the Hubble volume $\rho_\Lambda V_H=1/H=T_{\mathrm{GH}}\,S_{\mathrm{GH}}$ is the first law of de~Sitter, and the Komar/Landau--Lifshitz charge of the horizon $\kappa A/(8\pi\GN)$ gives the same $1/H$. The modulus $\muthree=\tfrac12$ is the ratio $T_{\mathrm{GH}}/T_{\mathrm{local}}$ of horizon and local temperatures. The gravitational field energy---absent from $\nabla^{A}T_{AB}=0$ (Proposition~\ref{prop:obs-matter})---is thus fixed, the energy conjugate to the tower entropy $S(\sigma)$ by the first law.
\ProofWitnessBlock{ds\_first\_law, energy\_is\_inv\_H, komar\_first\_law, mu3\_temp\_ratio, LL\_spatial\_ne\_zero}
\NumericalWitnessBlock{eq:obs-landauer}
\end{theorem}\begin{remark}[The G--$\Lambda$ interval: $\Lambda$ is determined, not selected]
\label{rmk:GLambda-interval}
The cosmological constant is not an input. The scale $\ell=1$ (from the
S-duality fixed point $\tau=i$, \eqref{eq:bridge-S}) and the dimension $d=4$
(from the arity $n=3$, \S\ref{ssec:arity}) are both derived, and they fix
\begin{equation}\label{eq:Lambda-from-curvature}
  \Lambda_5=-\frac{d(d-1)}{2\ell^{2}}=-6 ,
\end{equation}
with zero free parameters. The $6$ arrives by two independent routes---the curvature magnitude
$|\Lambda_5|=d(d-1)/2$ and the tower ceiling $\sigma_\Lambda=2n$ of \eqref{eq:interval-levels}---and
the routes agree only at the arity: $(n{+}1)n/2=2n$ forces $n=3$.
\ProofWitnessBlock{two\_routes\_to\_six}%
This realises precisely what
Witten~\cite{Witten01dS} identified as necessary: a theory
``that determines the value of the cosmological constant.''

\emph{The tower levels.} The activation levels of the three sectors are read
off the arity, not postulated:
\begin{equation}\label{eq:interval-levels}
  \sigma_G=n-1=2,\qquad \sigma_{\mathrm{EM}}=n=3,\qquad \sigma_\Lambda=2n=6 .
\end{equation}
Gravity activates at $\sigma_G=n-1$ (the spatial bulk $\E^{3}$ complete, one
dimension below the bulk); the electromagnetic sector sits at the midpoint
$\sigma_{\mathrm{EM}}=n$ (the Page point that halves the tower $[0,2n]$); and
$\Lambda$ activates at the ceiling $\sigma_\Lambda=2n$. The ceiling is
geometric: the torus $T^{2}_{\PCF}$ is a \emph{complex} space, so complex
dimension $n$ is real dimension $2n$---``complex dimension three'' is
$\R^{6}$, never $\R^{3}$. The tower runs over the full real space
$\sigma\in\{0,1,\dots,2n\}$, and its ceiling $\sigma_\Lambda=2n$ is that real
dimension.

\emph{The interval is the spacetime.} The gap between the gravity and
$\Lambda$ levels is exactly the dimension of spacetime,
\begin{equation}\label{eq:interval-gap}
  \sigma_\Lambda-\sigma_G=2n-(n-1)=n+1=4=\dim(M^{4}),
\end{equation}
and the interval carries the two holographic invariants of the microstate:
\begin{equation}\label{eq:interval-fractions}
  \frac{\sigma_{\mathrm{EM}}-\sigma_G}{\sigma_\Lambda-\sigma_G}
    =\frac{1}{n+1}=\frac14=|\Om|^{2}
  \qquad\text{and}\qquad
  \frac{\sigma_{\mathrm{EM}}-\sigma_G}{\sigma_\Lambda-\sigma_{\mathrm{EM}}}
    =\frac1n=\frac13=\|P\|^{2},
\end{equation}
the holographic area factor $\tfrac14=|\Om|^{2}$ and the past-component
weight $\tfrac13=\|P\|^{2}=F_\Om/c$.

\emph{Two independent routes to the same $6$.} The curvature magnitude
$|\Lambda_5|=d(d-1)/2=6$ (from $d=4$) and the tower ceiling
$\sigma_\Lambda=2n=6$ (from $n=3$, as the real dimension of the complex
torus) coincide by independent routes---5D curvature on one side, the
real dimension of the worldsheet torus on the other. Both are the entropy
saturation of the microstate read geometrically: $\sigma_\Lambda=2n=6$ is the
ceiling of the golden tower $S(\sigma)=\pi\golden^{\sigma}$ whose accumulated
bits are the horizon entropy of \eqref{eq:obs-identity} and, through
\eqref{eq:obs-einstein}, the field equations---the six-fold value is that
saturation, not a separate geometric fact. Together with $6=|S_{3}|=3!$ (the order of the symmetry
group, Lem.~\ref{lem:s3-orders}) and the exponent of the $G$--$\Lambda$ duality
$\golden^{-6}\golden^{+6}=1$, the value $\sigma_\Lambda=6$ is fixed from
four sides.

\emph{Scope.} This we propose as the structure Witten required: the de~Sitter
Hilbert-space dimension $N$ is integer, so $G\Lambda^{(n-2)/n}$ cannot vary
continuously and both $G$ and $\Lambda$ must be fixed by a more fundamental
theory~\cite{Witten01dS}; here $\GN=\tfrac12$ and $\Lambda\sim\golden^{-6}$, the
product fixed, $\golden^{-6}\golden^{+6}=1\iff G\Lambda\sim
M_{\mathrm{Pl}}^{-4}M_{\mathrm{Pl}}^{4}=1$ (PCF units). What is determined is this
dimensionless conjugacy. The translation to the observed
$\Lambda_{\mathrm{obs}}\sim10^{-122}M_{\mathrm{Pl}}^{4}$ is a matter of scale translation, not
derived here; and it is a \emph{single} open item rather than two. The level the observer occupies
is not an independent parameter: the tower fixes it from the de~Sitter entropy, and that entropy is
$1/\Lambda$, so
\begin{equation}\label{eq:sigma-obs}
  \sigma_{\mathrm{obs}}=\frac{\ln(S_{dS}/\pi)}{\ln\golden},\qquad
  S_{dS}=\frac{3\pi}{\GN\Lambda}\ \overset{\GN=1/2}{=}\ \frac{6\pi}{\Lambda}
  \quad\Longrightarrow\quad
  \sigma_{\mathrm{obs}}=\frac{\ln(6/\Lambda_{\mathrm{obs}})}{\ln\golden}.
\end{equation}
Deriving the scale translation therefore derives the level, and conversely; counting them as two
debts would be an accounting error. The scale $\sqrt M$ that sets the charged leptons (\S\ref{prop:obs-matter}) is therefore the tower reading of this de~Sitter entropy, and Theorem~\ref{thm:colour-from-M} then shows the colour sector hanging on the same scale: one scale surfaces as the string constant of \S\ref{sec:derivations} and as the entropy here, the dimensionless placement already fixed, the absolute value alone carrying the observer's $\Lambda$. What the tower does add is granularity: its step is exactly
$\golden$ ($\log_{10}\golden=0.209$ per level), so the absolute scale is constrained to a discrete
ladder rather than a continuum.
\NumericalWitnessBlock{eq:Lambda-from-curvature}
\ProofWitnessBlock{lambda\_magnitude\_eq\_level, sigma\_G\_val, sigma\_EM\_val, sigma\_Lambda\_val, threshold\_ordering, interval\_gap, interval\_gap\_general, em\_holographic\_fraction, em\_asymmetry, entropy\_at\_mu\_three, tower\_step\_ratio}
\end{remark}

\begin{proposition}[The level triple is unique over the integers]\label{prop:interval-uniqueness}
Consider integer triples $\sigma_G<\sigma_{\mathrm{EM}}<\sigma_\Lambda$ subject to the four
conditions already in place: the ceiling $\sigma_\Lambda=2n$, the gap
$\sigma_\Lambda-\sigma_G=n+1$ \eqref{eq:interval-gap}, and the two fractions of
\eqref{eq:interval-fractions} equal to the microstate invariants,
\begin{equation}\label{eq:interval-unique}
  \frac{\sigma_{\mathrm{EM}}-\sigma_G}{\sigma_\Lambda-\sigma_G}=|\Om|^{2}=\tfrac14,
  \qquad
  \frac{\sigma_{\mathrm{EM}}-\sigma_G}{\sigma_\Lambda-\sigma_{\mathrm{EM}}}=\|P\|^{2}=\tfrac13 .
\end{equation}
At the arity $n=3$ of \S\ref{ssec:arity} the four are satisfied by exactly one triple,
$(\sigma_G,\sigma_{\mathrm{EM}},\sigma_\Lambda)=(2,3,6)$, which is the assignment
\eqref{eq:interval-levels}. The selection is by the arity and not by fitting: the family
$(n-1,n,2n)$ satisfies the two fractions as $1/(n+1)$ and $1/n$ for \emph{every} $n$, and only
$n=3$ carries them onto $|\Om|^{2}=\tfrac14$ and $\|P\|^{2}=\tfrac13$ of
\eqref{eq:norms-derived}---at $n=2$ they are $\tfrac13,\tfrac12$ and at $n=4$ they are
$\tfrac15,\tfrac14$. Equivalently, $\sigma_\Lambda-\sigma_G=n+1$ equals $4=\dim(M^{4})$ only at
$n=3$. This is what Fig.~\ref{fig:alpha-uniqueness} displays --- panel (b) the constraint count,
panel (a) the spectral-angle surface of \S\ref{ssec:meta} evaluated on the resulting pair --- and it
is what lets
\S\ref{ssec:meta} read the spectral angle at the integer pair
$(\sigma_G,\sigma_{\mathrm{EM}})=(2,3)$ without circularity: the levels come from the arity, and
the angle is then evaluated on them. \ProofWitnessBlock{interval\_uniqueness, interval\_arity\_discriminates, interval\_gap\_only\_three}

\end{proposition}
\begin{proof}
Write $a=\sigma_{\mathrm{EM}}-\sigma_G>0$ and $b=\sigma_\Lambda-\sigma_{\mathrm{EM}}>0$, so
$\sigma_\Lambda-\sigma_G=a+b$. The second fraction of \eqref{eq:interval-unique} is $a/b=\tfrac13$,
i.e.\ $b=3a$; the first is then $a/(4a)=\tfrac14$ identically, so the two fractions are one
condition. The gap gives $a+b=4a=n+1$, whence $a=1$ at $n=3$ and $b=3$; the ceiling gives
$\sigma_\Lambda=6$, hence $\sigma_G=6-4=2$ and $\sigma_{\mathrm{EM}}=3$. Uniqueness is that each
step determined its unknown. For the discrimination, $(n-1,n,2n)$ gives $a=1$, $b=n$, $a+b=n+1$, so
the fractions are $1/(n+1)$ and $1/n$; these equal $\tfrac14$ and $\tfrac13$ iff $n=3$.
\end{proof}

\begin{figure}[!htbp]\centering
\includegraphics[height=0.82\textheight,keepaspectratio]{fig1_alphas_uniqueness.pdf}
\caption{(a) The spectral-angle surface $\sin\alpha(\sigma_1)\cos\alpha(\sigma_2)$ of
\eqref{eq:spectral-surface}; $\tan\alpha(\sigma{+}1)/\tan\alpha(\sigma)=\golden$
(\eqref{eq:spectral-angle}). The marked point $(2,3)$ gives
$\sin\alpha(2)\cos(3)=0.1181\approx\sin\alpha_{\rm GUT}$.
(b) Constraint count for each triple $(\sigma_G,\sigma_{\mathrm{EM}},\sigma_\Lambda)$:
$\sigma_\Lambda=2n$, $\sigma_\Lambda-\sigma_G=n{+}1$,
$(\sigma_{\mathrm{EM}}{-}\sigma_G)/(\sigma_\Lambda{-}\sigma_G)=|\Om|^2$,
$(\sigma_{\mathrm{EM}}{-}\sigma_G)/(\sigma_\Lambda{-}\sigma_{\mathrm{EM}})=\|P\|^2$.
Only $(2,3,6)$ satisfies all four (boxed). No fitted quantity enters.}
\label{fig:alpha-uniqueness}
\end{figure}

\subsection{Why this is a single demonstration}\label{ssec:spine4}

The path went through the three types: it opened in the non-local type~III field, made the
type~I particle emerge, turned it into the type~II observer, interwove I and II in the
tower, and closed in the type~III spacetime. These are not five demonstrations glued
together but one, because the same object of modulus $\tfrac12$ is the protagonist of all,
and each transition is a proved implication: entropy orients the flow and yields the
particle, the Wick rotation gives the cones and the observer, accumulation overflows and
builds the tower, entropy sources gravity.

The commutative diagram supports this unity with precise structure. The twenty-four
equations form a directed acyclic graph with two roots---the projection $\Pi$
(\eqref{eq:obs-interface}) and the observer's objective record
(\eqref{eq:obs-redundancy})---from which all else derives, with $36$ demonstrative edges.
A topological order exists---the equations can be stated so that none depends on a later
one---which guarantees the section reads without forward reference. Eight nodes are apices,
where two or more independent chains converge; the weld (\eqref{eq:obs-weld}), with four
routes, is the maximal one and the only triple point I${\cap}$II${\cap}$III. The circle
closes literally: \eqref{eq:obs-einstein} returns to \eqref{eq:obs-interface}---the
non-local gravity returns to the same $\Pi$ and $\tfrac12$ of the opening, III${\to}$III.
The spacetime at the end is non-local at every point, type~III, as the field at the start,
and its Newton constant is the same $\tfrac12$ of the object one began with. One enters by
any door and reaches the same object: that is what it means for the diagram to commute.
The section does not prove six things; it proves that an observer of modulus $\tfrac12$,
which cannot be separated into parts, is at once the particle, the string, the tower of
information and the curved de~Sitter universe that contains it, and that this universe has
the $\Lambda$ the theory was bound to determine, with the observer inside. The debt with
which the section opened is settled by the same object that raised it.

\medskip\noindent\textbf{The demonstrative line.} Each step uses only the previous ones.

\begin{figure}[htb]\centering
\resizebox{\textwidth}{!}{%
\begin{tikzpicture}[
  font=\footnotesize,
  >={Stealth[length=4pt]},
  nd/.style={inner sep=2pt, outer sep=1pt},
  ap/.style={draw, rounded corners=2pt, fill=black!5, inner sep=3pt},
  every path/.style={->, draw=black!72, line width=0.45pt}]

% ===== 4.2 field + particle (x=0) =====
\node[nd] (int) at (0,5.4)  {$\Pi\colon\E^3{\to}\C$};
\node[nd] (spin) at (0,4.2) {$C{+}P{+}F$};
\node[nd] (half) at (0,3.0) {$\|P\|\|C\|\|F\|{=}\tfrac12$};

% ===== 4.3 observer (x=3.4) =====
\node[nd] (thr)  at (3.4,2.0) {$f_{\rm crit}{=}\tfrac12$};
\node[nd] (cram) at (3.4,0.9) {$\Fmax^{-1}{=}\tfrac14$};
\node[nd] (fish) at (3.4,-0.2){$\tau_F{=}\tau_D{\Leftrightarrow}f{=}\tfrac12$};
\node[nd] (red)  at (3.4,3.2) {$R_\delta{\sim}N$};
\node[ap] (acc)  at (3.4,4.4) {$F_\theta{\to}\Fmax{=}4$};
\node[nd] (cert) at (3.4,5.6) {$\varepsilon_0 M{=}\pi$};

% ===== 4.4 tower (x=6.8) =====
\node[nd] (fix)    at (6.8,6.2) {$\beta_g{=}0$};
\node[nd] (throat) at (6.8,4.8) {$z{=}\golden^\sigma,S{=}\pi\golden^\sigma$};
\node[nd] (swp)    at (6.8,3.4) {$\tfrac{|\partial_\sigma V|}{V}{=}\ln\golden$};
\node[ap] (weld)   at (6.8,2.0) {$\tau_F{=}\tau{=}M\golden^{-\sigma}$};
\node[ap] (id)     at (6.8,0.6) {$F_\Om\Nmodes{=}\pi\golden^\sigma$};

% ===== 4.4 gravity (x=10.4) =====
\node[ap] (land) at (10.4,5.6) {$E/\text{bit}{=}1/M$};
\node[nd] (jac)  at (10.4,4.4) {$\delta Q{=}T\delta S$};
\node[nd] (ein)  at (10.4,3.2) {$R_{AB}{=}{-}4g_{AB}$};
\node[ap] (ets)  at (10.4,2.0) {$ds^2{=}d\vec x^2{-}c^2dt^2$};
\node[ap] (mat)  at (10.4,0.6) {$\text{matter}{=}\lfloor S\rfloor$};

% ===== 4.5 Lambda + interval (x=13.8) =====
\node[nd] (lam)  at (13.8,4.2) {$\Lambda_5{=}{-}6$};
\node[nd] (lev)  at (13.8,3.0) {$\sigma_G,\sigma_{\rm EM},\sigma_\Lambda$};
\node[nd] (gap)  at (13.8,1.8) {$\sigma_\Lambda{-}\sigma_G{=}4$};
\node[ap] (frac) at (13.8,0.6) {$\tfrac1{n+1}{=}\tfrac14{=}|\Om|^2$};
\node[ap] (uniq) at (13.8,-0.7) {$(\sigma_G,\sigma_{\rm EM},\sigma_\Lambda){=}(2,3,6)$ unique};

% ===== edges (columna a columna, como spine3) =====
\draw (int) -- (half);
\draw (int) -- (spin);
\draw (spin) -- (acc);
\draw (half) to[bend left=5] (thr);
\draw (thr) -- (cram);
\draw (thr) -- (fish);
\draw (cram) -- (acc);
\draw (red) -- (acc);
\draw (acc) -- (cert);
\draw (cert) -- (fix);
\draw (cert) -- (throat);
\draw (throat) -- (swp);
\draw (fix) -- (weld);
\draw (acc) to[bend right=8] (weld);
\draw (throat) -- (weld);
\draw (cert) to[bend left=8] (weld);
\draw (weld) -- (id);
\draw (throat) to[bend left=6] (id);
\draw (acc) to[bend right=14] (id);
\draw (id) to[bend right=8] (land);
\draw (cert) to[bend left=20] (land);
\draw (land) -- (jac);
\draw (jac) -- (ein);
\draw (throat) to[bend left=4] (mat);
\draw (id) -- (mat);
\draw (ein) -- (mat);
\draw (throat) to[bend right=4] (ets);
\draw (half) to[bend right=10] (ets);
\draw (weld) -- (ets);
\draw (ein) -- (lam);
\draw (lam) -- (lev);
\draw (lev) -- (gap);
\draw (lev) -- (frac);
\draw (half) to[out=-30,in=180,looseness=0.5] (frac);
\draw (frac) -- (uniq);
\draw (gap) to[bend left=10] (uniq);
% cierre III->III
\draw[dashed, draw=red!55] (ein) to[out=-90,in=-90,looseness=0.5]
  node[below,font=\scriptsize]{III$\to$III} (int);

\end{tikzpicture}
}
\caption{Commutative diagram of \S\ref{sec:implications}: the twenty-four equations
of \S4 as nodes, arranged by subsection (the non-local field and particle,
\S\ref{ssec:hinge}; the observer, \S\ref{ssec:observer}; the tower and gravity,
\S\ref{ssec:accum}), with the thirty-six demonstrative arrows. Shaded nodes are the
apices, where two or more independent chains converge on the same value; the weld
$\tau_F{=}\tau{=}M\golden^{-\sigma}$ has four incoming routes and is the maximal apex.
The dashed return edge is the closure $R_{AB}{=}-4g_{AB}\to\Pi$ (type~III to type~III):
the non-local gravity returns to the same projection $\Pi$ and modulus $\tfrac12$ of the
opening.}
\label{fig:spine4}
\end{figure}

\begin{description}
  \item[\S4.1 what is proved, and against what.] Holography in de~Sitter requires a theory that
        determines $\Lambda$ and describes the observer from within (Witten). The yardstick is the
        von~Neumann classification into three types---I discrete (particle, definite spectrum), II
        observation (trace, crossed product), III local field (worldsheet, horizon,
        non-locality)---and what is proved is the \emph{signature} of each, not the technical name.
        The criteria come from three distinct sources, each tagged where used: Witten~2001
        ($\Lambda$, finiteness $S=\ln N$, the isometry obstruction, the pairing, the CPT map, the
        two conjectures), CLPW~2022 (the observer described by the theory, the maximally mixed
        vacuum, the crossed product), and Gibbons--Hawking (the thermal static patch). The thesis:
        the three types are realised on one object, the microstate $\Om$ with $|\Om|=\tfrac12$, as
        three threads crossing the whole section.
  \item[\S4.2 the field and the particle.] The path opens in the non-local type~III field:
        the worldsheet carries the interface $\Pi\colon\E^3\to\C$ (\eqref{eq:obs-interface}),
        whose three components are the spin-star $C$-$P$-$F$ (\eqref{eq:obs-spinstar}). The whole
        geometry is one tensor---the warp $z=\golden^\sigma$ generates the curvature
        $R_{AB}={-}4g_{AB}$ (App.~\ref{app:curvature})---and from that same geometry, by its
        tension, the discrete type~I particle emerges. Geometry and particle are the same object.
  \item[\S4.3 the particle as observer.] An internal Wick rotation turns the particle into an
        observer: the centre becomes the time axis, past and future the two cones. The observer
        accumulates---Fisher grows to the Cram\'er--Rao ceiling $\Fmax^{-1}=\tfrac14$
        (\eqref{eq:obs-cramerrao}, \eqref{eq:obs-fishertime})---and sees the exact half
        $\|P\|\|C\|\|F\|=\tfrac12=|\Om|$ (\eqref{eq:obs-half}), with redundancy $R_\delta\sim N$
        (\eqref{eq:obs-redundancy}) and accumulation to $\Fmax=4$ (\eqref{eq:obs-accum}). Three
        independent routes converge on $\tfrac12$ (first apex), the objectivity threshold
        $f_{\mathrm{crit}}=\tfrac12$ of Kiely (\eqref{eq:obs-threshold}) and the cell capacity
        $\epsz\Mpcf=\pi$ (\eqref{eq:obs-certainty}).
  \item[\S4.4 the tower, and curved spacetime.] The accumulation builds the tower (types~I and~II
        interwoven): the cell overflows and the throat opens, $z=\golden^\sigma$,
        $S(\sigma)=\pi\golden^\sigma$ (\eqref{eq:obs-throat}, \eqref{eq:obs-swampland}), at the UV
        fixed point $\beta_g=0\Leftrightarrow\epsz\Mpcf=\pi$ (\eqref{eq:obs-fixedpoint}). The two
        accumulations---of time and of scale---are one, the weld (\eqref{eq:obs-weld}), giving the
        identity $F_\Om\Nmodes=S(\sigma)=\SBH=S_{\text{Jacobson}}$ (\eqref{eq:obs-identity}). Then
        curved spacetime (type~III again): Landauer (\eqref{eq:obs-landauer}) $\to$ Jacobson
        (\eqref{eq:obs-jacobson}) $\to$ Einstein (\eqref{eq:obs-einstein}), with matter the
        saturation entropy (\eqref{eq:obs-matter}). The internal Wick rotation gives the de~Sitter
        metric (\eqref{eq:ets-metric}), $\Lambda_5={-}6$ (\eqref{eq:Lambda-from-curvature}), and the
        level interval $\sigma_G{=}2,\sigma_{\mathrm{EM}}{=}3,\sigma_\Lambda{=}6$
        (\eqref{eq:interval-levels}) with $\sigma_\Lambda{-}\sigma_G{=}n{+}1{=}4$
        (\eqref{eq:interval-gap}) and the fractions $\tfrac1{n+1}{=}\tfrac14{=}|\Om|^2$
        (\eqref{eq:interval-fractions}); the $G$--$\Lambda$ duality and
        $\sin^2\theta_W|_{\rm GUT}{=}\tfrac38$ (Prop.~\ref{prop:obs-matter}) give the gauge factors
        $\mathrm{SU}(3){\times}\mathrm{SU}(2){\times}\mathrm{U}(1)$ (Fig.~\ref{fig:torus-gauge}).
\end{description}
\noindent The observer does not select; it builds without alternative --- participatory yet
forced --- consonant with Kiely's quantum Darwinism. The whole section is one commutative
diagram (Fig.~\ref{fig:spine4}): the twenty-four equations as nodes, the demonstrative
arrows between them, and the apices where independent routes converge.

% ===================================================================
\FloatBarrier
\section{Implications for holography regarding dynamics and operators}\label{sec:dof}

\S\ref*{sec:derivations}--\S\ref*{sec:implications} showed that accumulated entropy becomes
gravity and geometry. This section proves the companion statement: the same accumulated entropy
becomes degrees of freedom. The reading is not a second theory but a second face of the first
one; the object is the microstate $\Om$ with $\lvert\Om\rvert=\tfrac12$, and the yardstick is the
positive Grassmannian, on which one datum is at once an amplitude, a fluid, a particle and a
state. \S\ref*{ssec:adscft} derived the whole of AdS/CFT from the single microstate but deferred its
realisation as a self-adjoint operator---the absence this paper opened with, and why
low-dimensional holography recovers its bulk as an ensemble average. This section constructs
our proposal for it. The dynamics is forced: the eigenvalue triangle is the $A_2$ root system,
$A_2$ is the root system of $\su(3)$, and a degree of freedom is a point of the positive
Grassmannian, so what the object already carries is a gauge theory. Nor is an operator
enough---to stand as a boundary dual it must be gapped, since it is the gap that makes the
Osterwalder--Schrader reconstruction return a field theory. The section computes the tower
Hamiltonian, the strong-coupling gap $\Delta=\tfrac23g^{2}$, and the reconstruction they
support; the isometry of \eqref{eq:isometry} then intertwines that Hamiltonian with the tower
entropy level by level (\S\ref{ssec:operators}). Each step uses only earlier ones, and two
invariants already fixed in \S\ref*{sec:methods}--\S\ref*{sec:derivations}
--- the tower/Euler product $\zeta$ and the modulus $\tfrac12$ --- thread every movement. Tiers are
used throughout as in \S\ref{sec:methods}.

\begin{table}[htb]\centering\footnotesize
\caption{Dictionary: framework object $\leftrightarrow$ standard object (cf.\ the D1--D11 dictionary of \S\ref*{sec:derivations}).}\label{tab:dict}
\begin{tabular}{lll}\hline
framework & standard & \S\\\hline
$\lvert\Om\rvert=\tfrac12$ & modulus / spin $\tfrac12$ & \ref{ssec:grass}\\
$\rho=P/k$ & density matrix & \ref{ssec:grass}\\
tower $\lfloor\pi\golden^\sigma\rfloor$ & Regge trajectory & \ref{ssec:amp}\\
$\zeta=\prod_p$ & Euler product / OPE & \ref{ssec:amp}\\
$H\lvert\sigma\rangle=m_0\golden^\sigma$ & scale operator (KK side) & \ref{ssec:operators}\\
$C_2(n)$ & colour charge & \ref{ssec:operators}\\
$T=e^{-aH}$ & transfer operator & \ref{ssec:os}\\
$V^\dagger V=\mathbf1$ & holographic isometry & \ref{ssec:os}\\
$\Theta$ & CPT reflection & \ref{ssec:os}\\
$\xi(s)=\xi(1-s)$ & S-duality & \ref{ssec:zeta-gap}\\\hline
\end{tabular}\end{table}

\subsection{The entropy\texorpdfstring{$\to$}{->}DOF law}\label{ssec:law}
\S\ref*{ssec:accum} let the observer's information overflow into the tower $S(\sigma)=\pi\golden^{\sigma}$
and, by Jacobson, into curved spacetime. Read the same overflow not as area but as \emph{count}:
each cell that the tower opens is a degree of freedom, and their number grows at the rate the
tower grows. Bianconi's entropic gravity~\cite{Bianconi}, in which a relative entropy of the metric sources the field, is the classical and continuous limit of this count; measured against CW6 it reduces on
five object-level squares.

\begin{proposition}[Five reduction squares]\label{prop:bianconi}
Under the reduction $R$ (divide by $4$, $\mathrm{Berry}\to0$, discrete$\to$continuous), the CW6
data map to Bianconi's:
\begin{align*}
\text{(D1)}\ & \mathrm{QFI}_{\mathrm{CW6}}/4=M_{\mathrm{Bianconi}}, & \text{(Fisher; the }4\text{ from }\mu=\tfrac12)\\
\text{(D2)}\ & \Nmodes=\lfloor\pi\golden^{\sigma}\rfloor\ \xrightarrow{R}\ S=\pi\golden^{\sigma}, & (|\ln\Nmodes-\ln S|\to1.4\times10^{-7})\\
\text{(D3)}\ & S=kA,\ k=\tfrac1{4G_N}=\mu=\tfrac12, & \text{(area law)}\\
\text{(D4)}\ & \Lambda_G\approx\tfrac1{4\beta}\varepsilon^2\ \Rightarrow\ \beta=\tfrac{\ln2}{8}, & \text{($\Lambda$ scale, cf.\ }\Lambda_5=-6\text{, \S\ref{ssec:adscft})}\\
\text{(D5)}\ & \{\tfrac12\omega^{k}\}_{k=0}^{2}\ \leftrightarrow\ 0\oplus1\oplus2, & \text{(triad $\to$ Dirac--K\"ahler)}
\end{align*}
The five squares commute, and the relation they express is a \emph{limit} one: Bianconi's
continuous entropic gravity is the continuum limit of the discrete accumulation of
\S\ref{ssec:accum}---D2 is exactly that limit, the floor removed. What the present framework adds
to her axis is the quantity it leaves free, the \emph{rate}: the growth is golden, $\golden$ per
level, the Fibonacci quantum dimension (Corollary~\ref{cor:golden-growth}). \NumericalWitnessBlock{cor:golden-growth}
D3 (\S\ref*{ssec:construction}). \ProofWitnessBlock{binary\_entropy\_half}
\end{proposition}
\begin{proof}
D3 is \eqref{eq:bridge-BH}. For D1, the CW6 quantum Fisher information at the operating point
$\mu=\tfrac12$ is $F=4(M-|\mathrm{Berry}|^2)$; setting $\mathrm{Berry}\to0$ and dividing by $4$
gives $M$. D2 is the discrete-to-continuous limit of \eqref{eq:tower-modes}, the floor removed.
D4 matches the coefficient of the small-parameter expansion of $\Lambda_G$ to CW6's
$1-H\approx(2/\ln2)\delta^2$, fixing $\beta$; here $\varepsilon$ and $\delta$ are identified as the
same dimensionless deviation from the maximally mixed point ($\hat G=\mathbf 1$ on that side,
$p=\tfrac12$ on this one), which is what makes the two quadratics comparable. D5 identifies the three eigenvalues of
Proposition~\ref{prop:pcf-norms} with the three form-grades; on that side the truncation at
$2$-forms is stated as a simplification admitting extension to higher forms, so the identification
is tighter from the present side than from hers. The two horizontal (dynamical) arrows
are distinct routes to the same conclusion---both derive Einstein-type field equations from
entropy: Bianconi's by varying a quantum relative entropy of the metric, the present one
thermodynamically through $\delta Q=T\,\delta S$ on local horizons (\S\ref{ssec:accum},
Jacobson). The correspondence asserted here is the commutation of the vertical maps together with
the limit relation of D2; it does not assert that the two derivations are the same map. The
reduction is also silent on one axis: that programme derives the area law \emph{without} the
holographic principle, whereas the present one derives it with the bulk--boundary isometry
\eqref{eq:isometry}. Neither subsumes the other there; on whether holography is required the two
are complementary.
\end{proof}

\begin{corollary}[Golden growth of degrees of freedom]\label{cor:golden-growth}
The number of degrees of freedom at scale $\sigma$ grows as $\golden^{\sigma}$: the count inherits
the tower's rate, the successive ratio $\Nmodes(\sigma)/\Nmodes(\sigma-1)\to\golden$ of
Fig.~\ref{fig:tower-modes}. \ProofWitnessBlock{S\_tower, tower\_step\_ratio}

\end{corollary}

From here the degrees of freedom, not the geometry, are followed; the next movement gives them a
home.

\subsection{The positive Grassmannian}\label{ssec:grass}
A degree of freedom is not a bare number but a point of a positive geometry~\cite{Postnikov,BengtssonZyczkowski}, and that geometry is
already a state.

\begin{proposition}[Grassmannian point is a density matrix]\label{prop:rho}
A point of $\Grpos(k,n)$ with full-rank frame $C$ defines the orthogonal projector
$P=C^{\mathsf T}(CC^{\mathsf T})^{-1}C$, with $P^2=P$ and $P^{\mathsf T}=P$, and the density
matrix $\rho=P/k$, positive semidefinite with $\operatorname{tr}\rho=1$. For $k=2$ its nonzero
eigenvalues equal $\tfrac12=\lvert\Om\rvert$. \ProofWitnessBlock{projector\_idem, rho\_is\_state, rho\_k2\_eq\_half\_projector, projector\_trace\_eq\_rank}

\end{proposition}

\begin{proposition}[The two positivities are carried by one projector]\label{prop:two-pos}
For a full-rank frame $C$ the projector supplies at once the quantum positivity of the state and the
total positivity of the geometry:
\begin{equation}\label{eq:two-pos}
  \rho=P/k\succeq0,\qquad\text{and}\qquad C\in\Grpos \Longrightarrow \Delta_I(C)\ge0
  \ \text{ for every ordered } I .
\end{equation}
The first is automatic from idempotence and holds for \emph{any} point of the Grassmannian; the
second is the total positivity of the Pl\"ucker minors and holds only on $\Grpos$. They are therefore
not the same condition: the second is strictly stronger, and it is the \emph{refinement}---one datum
carrying an operator-level positivity and an amplitude-level positivity at different levels---that
lets the state reading and the amplitude reading be taken on one object. \ProofWitnessBlock{two\_positivities\_distinct, rho\_k2\_eq\_half\_projector}

\end{proposition}
\begin{proof}
$P^2=C^{\mathsf T}(CC^{\mathsf T})^{-1}(CC^{\mathsf T})(CC^{\mathsf T})^{-1}C=P$; symmetry follows
from $(CC^{\mathsf T})^{\mathsf T}=CC^{\mathsf T}$. A projector is positive semidefinite with
eigenvalues in $\{0,1\}$; dividing by $k$ gives $\rho\succeq0$ and, since $\operatorname{tr}P=\mathrm{rank}\,P=k$,
$\operatorname{tr}\rho=1$. For $k=2$ the nonzero eigenvalues are $1/2$.
\end{proof}

\begin{theorem}[One datum, four faces]\label{thm:faces}
Let $C$ be a full-rank $k\times n$ frame of a point of $\Grpos(k,n)$ and
$P=C^{\mathsf T}(CC^{\mathsf T})^{-1}C$. Then $P$ depends only on the point and not on the frame:
for every invertible $g$,
\begin{equation}\label{eq:frame-invariance}
  P(gC)=P(C),
\end{equation}
and the four readings of \S\ref{sec:dof} are four functions of that one point, all factoring
through $P$:
\begin{equation}\label{eq:four-faces}
  \underbrace{P^{2}=P,\ P^{\mathsf T}=P}_{\text{geometry}},\qquad
  \underbrace{\operatorname{tr}P=k}_{\text{degrees of freedom}},\qquad
  \underbrace{\rho=P/k,\ \operatorname{tr}\rho=1}_{\text{state}},\qquad
  \underbrace{\text{maximal minors}\ \ge 0}_{\text{fluid}} ,
\end{equation}
the last for $C$ totally positive (Proposition~\ref{prop:fluid} below). The canonical form of the
positive geometry is a function of the positroid cell, hence of the same point
(\S\ref{ssec:amp}). \ProofWitness{projector\_frame\_invariant}
%
\ProofWitnessBlock{projector\_idem, projector\_trace\_eq\_rank, rho\_is\_state, positroid\_from\_kp}%
\end{theorem}
\begin{proof}
For \eqref{eq:frame-invariance}, $(gC)^{\mathsf T}\bigl(gC(gC)^{\mathsf T}\bigr)^{-1}(gC)
=C^{\mathsf T}g^{\mathsf T}\bigl(g(CC^{\mathsf T})g^{\mathsf T}\bigr)^{-1}gC
=C^{\mathsf T}g^{\mathsf T}(g^{\mathsf T})^{-1}(CC^{\mathsf T})^{-1}g^{-1}gC
=C^{\mathsf T}(CC^{\mathsf T})^{-1}C$, using $(gAg^{\mathsf T})^{-1}=(g^{\mathsf T})^{-1}A^{-1}g^{-1}$.
The four clauses of \eqref{eq:four-faces} are Proposition~\ref{prop:rho} and its proof
($P^2=P$, $P^{\mathsf T}=P$, $\operatorname{tr}P=k$, $\operatorname{tr}\rho=1$) together with
Proposition~\ref{prop:fluid} below for the minors. What \eqref{eq:four-faces} asserts is not that four
numbers coincide but that four constructions share an argument: each is a function of the point,
and by \eqref{eq:frame-invariance} none of them can distinguish two frames of the same point.
That is the precise content of $\text{amplitude}=\text{fluid}=\text{DOF}=\text{state}$.
\end{proof}

\subsection{The ladder \texorpdfstring{$A_2\to E_8$}{A2->E8}}\label{ssec:ladder}
The seed of the ladder is the triad already in hand.

\begin{proposition}[The $A_2$ seed]\label{prop:a2}
The eigenvalue triad $\lambda_k=\tfrac12\omega^{k}$ (Proposition~\ref{prop:pcf-norms}) is at once
the Eisenstein lattice $\Z[\omega]=A_2$, the root system of $\su(3)$, the finite cluster type of
$\mathrm{Gr}(3,5)$, the Dirac--K\"ahler grading $0\oplus1\oplus2$, and three bits ($c=3$, the central charge the
three eigenvalues give by Brown--Henneaux, \S\ref{ssec:adscft}; the count three itself is forced by
the arity uniqueness of \S\ref{ssec:arity}).
\ProofWitnessBlock{eisenstein, omega\_eigenvalues, su3\_dim}
\end{proposition}
\begin{proof}
$\omega=e^{2\pi i/3}$ generates $\Z[\omega]$, whose minimal vectors form the six roots of $A_2$
(\S\ref*{app:gauge}); $\dim\su(3)=3^2-1=8$; Scott's classification gives $\mathrm{Gr}(3,5)=A_2$;
the three eigenvalues are the three grades; and $\lfloor\pi\rfloor=3$ fixes $c=3$
(\S\ref*{app:superpoint}).
\end{proof}

\begin{proposition}[The finite-type ladder]\label{prop:ladder}
By Scott's classification~\cite{Scott,FZ} $\mathrm{Gr}(3,6)=D_4$, $\mathrm{Gr}(3,7)=E_6$, $\mathrm{Gr}(3,8)=E_8$;
the positive-root counts $3<12<36<120$ obey $\text{kissing}=2\cdot\text{roots}$ and dimensions $(2,4,6,8)$~\cite{ConwaySloane,Viazovska}, with $A_2$ the seed of each rung. The gauge-algebra dimensions
follow from the same data by $\dim\mathfrak g=\text{kissing}+\text{rank}$, the Frenkel--Kac--Segal
enhancement count at the self-dual point: $8=6+2$ for $\mathfrak{su}(3)$, $28=24+4$ for
$\mathfrak{so}(8)$, $78=72+6$ for $\mathfrak e_6$ and $248=240+8$ for $\mathfrak e_8$---so the
combinatorial ladder is, rung by rung, the ladder of gauge groups (\S\ref{app:kk}).
The dimension count is certified by the same data.
\ProofWitnessBlock{ladder\_from\_scott, scott\_kissing\_eq\_two\_posroots, kissing\_eq\_two\_posroots, dims\_even\_ladder, ladder\_dim\_eq\_kissing\_plus\_rank, fks\_ladder\_dims}

\end{proposition}

\subsection{First face: amplitude}\label{ssec:amp}
The amplitude is where $\zeta$ enters, and it enters already proved in \S\ref*{ssec:p-tower} and completed over all places in \S\ref*{ssec:zeta} (Theorem~\ref{thm:places}).

\begin{proposition}[Veneziano, Regge, Euler]\label{prop:veneziano}
The Veneziano amplitude~\cite{Veneziano,ArkaniHamedTrnka,GGSA} $A_4=\Gamma(-\alpha's)\Gamma(-\alpha'u)/\Gamma(1-\alpha'(s+u))$ has simple poles at
$\alpha's=n$, the Regge tower; the pole positions index $\sum_{n\ge1}n^{-s}=\zeta(s)=\prod_p(1-p^{-s})^{-1}$
(\S\ref*{ssec:p-tower}, \eqref{eq:euler-product}). \ProofWitnessBlock{veneziano\_eq\_gamma\_ratio, regge\_tower\_is\_euler\_product}

\end{proposition}

\begin{proposition}[Field-theory limit; the $\pi/\golden$ split]\label{prop:ftlimit}
$\displaystyle\lim_{\alpha'\to0^{+}}\alpha'\,B(\alpha's,\alpha't)=\tfrac1s+\tfrac1t=\Om(X)|_{n=4}$,
and the amplitude admits two computations meeting at $\lvert\Om\rvert=\tfrac12$ and $\zeta$: a
$2$D path integral (the $\pi$ side, $\Gamma(\tfrac12)=\sqrt\pi$) and a $3$D positive-geometry
canonical form (the $\golden$ side). \ProofWitnessBlock{ft\_identity, ft\_limit, pentagon\_u\_eq\_golden}
\end{proposition}

\subsection{Second face: particle, graviton, matter}\label{ssec:particle}
The spectrum of the object is a particle; its spin is the modulus, and its interactions carry the
gauge charges the torus already gave.

\begin{proposition}[Fermion and Standard-Model content]\label{prop:sm}
The self-conjugate midpoint of $x\mapsto1-x$ is at once $\lvert\Om\rvert$ and the fermion spin
$\tfrac12$ (\S\ref*{ssec:hinge}). The one-loop $\beta$-coefficients are $b=(\tfrac{33}5,1,-3)$;
each generation is anomaly-free, $\sum Y=0$; and $\sin^2\theta_W|_{\rm GUT}=N(0)/N(2)=\tfrac38$, running to $0.231$ at $M_Z$ (\S\ref*{app:gauge}).
\ProofWitnessBlock{beta\_is\_mssm, sm\_anomaly\_free, weinberg\_angle\_gut}
\end{proposition}

\begin{theorem}[The graviton is the entropy-response quantum]\label{thm:graviton}
About the ETS background---intrinsically flat, with de~Sitter its embedded hyperboloid
(\S\ref{app:embedding}, \S\ref{app:curvature})---the transverse--traceless perturbation $h_{\mu\nu}$
obeys $\square h_{\mu\nu}=8\pi G_N T_{\mu\nu}$ (massless spin-2) and coincides with the double copy~\cite{BCJ} $\text{graviton}=\text{YM}^2$ (spin $1+1=2$).
Read through \S\ref*{app:einstein} ($\delta Q=T\delta S\Rightarrow$ Einstein), it is the quantum
of the entropy response: the metric perturbation conjugate to the tower's Fisher information.
Quantitatively: the response rate is $S'(\sigma)=(\ln\golden)\,S(\sigma)$, the price per bit is the
constant temperature $\varepsilon(\sigma)/S(\sigma)=\varepsilon_0/\pi=1/\Mpcf$, and the Fisher clock
is the dynamical clock exactly at the modulus, $\tau_D/\sqrt{2f}=\tau_D\iff f=\tfrac12$
\eqref{eq:obs-fishertime}---so the $\delta S$ that sources $h_{\mu\nu}$ counts Fisher bits on the
tower's own clock.
; the double copy~\cite{BCJ} (Bern--Carrasco--Johansson) \ProofWitnessBlock{fisher\_time\_lock, entropy\_response\_per\_bit, entropy\_response\_has\_deriv\_at}

\end{theorem}
\begin{proof}
The linearised Einstein equation about the ETS background gives the wave equation for the
transverse--traceless part, sourced by $T_{\mu\nu}$ (\S\ref*{app:einstein}); masslessness is the
absence of a potential term. The colour--kinematics duality replaces each colour factor by a
second kinematic numerator, so the gravity numerator is the square of the gauge one, spin
$1+1=2$. Since the Einstein equation itself is $\delta Q=T\delta S$, the field $h_{\mu\nu}$ is the
first-order response of the geometry to a change in the accumulated entropy, i.e.\ to the tower's
Fisher information of \S\ref{ssec:law}. The three quantitative pieces close the reading: the rate,
the constant price per bit, and the clock lock at $f=\mu_3$ are each proved above, and the wave
operator itself is exercised numerically (any profile $h(t{-}z)$ is annihilated, a $w$-independent
mode feels no potential).
\end{proof}

\subsection{Supersymmetry as ER=EPR bridges}\label{ssec:susy}
The doubling that supersymmetry~\cite{Martin} performs is, here, not new particles but the entanglement bridges~\cite{MaldacenaSusskind13} already proved to form a group.

\begin{theorem}[Running without superpartners]\label{thm:susy}
The ER=EPR bridges $T(\sigma_1,\sigma_2)=D(\sigma_1)/D(\sigma_2)$ form a groupoid
(\S\ref*{ssec:hinge}) and carry the gauge representations of the fields they bridge. Their
contribution to the running satisfies
\begin{equation}\label{eq:susy-bridges}
  b_{\mathrm{MSSM}}=b_{\mathrm{SM}}+\Delta b_{\mathrm{bridges}},
\end{equation}
reproducing the MSSM one-loop coefficients with no superpartner particles. \ProofWitnessBlock{mssm\_eq\_sm\_plus\_bridges, bridge\_cocycle}

\end{theorem}
\begin{proof}
The groupoid laws (inverse, composition, cocycle) are \S\ref*{ssec:hinge}. A bridge of a field in
representation $R$ lies in $R$, so it contributes to the $\beta$-function with the sign of a
superpartner and the Casimir of $R$; summing the bridge contributions over the SM content gives
$\Delta b_{\mathrm{bridges}}$, and \eqref{eq:susy-bridges} is the arithmetic identity that this sum
carries $b_{\mathrm{SM}}$ to $b_{\mathrm{MSSM}}$.
\end{proof}

\begin{proposition}[State is geometry; the bridge splits the identity]\label{prop:er-epr-id}
For a frame $C$ with $k\neq0$ the density matrix rescaled by $k$ \emph{is} the projector, and any
idempotent splits the identity into orthogonal halves:
\begin{equation}\label{eq:er-epr-id}
  k\,\rho = P,\qquad P+(\mathbf1-P)=\mathbf1,\qquad P(\mathbf1-P)=0 .
\end{equation}
The first identity is $\mathrm{ER}=\mathrm{EPR}$ in the present language---the entangled state and
the geometric projector are the same operator up to the rank---and the second exhibits the bridge as
the decomposition of the identity that the groupoid of \S\ref{ssec:hinge} composes. \ProofWitnessBlock{er\_epr\_state\_eq\_geometry, er\_epr\_bridge}

\end{proposition}

\begin{corollary}[The bridges reproduce the supersymmetric coefficients]\label{cor:bridges-mssm}
Summing the bridge contributions over the Standard-Model content gives, componentwise,
\begin{equation}\label{eq:bridges-mssm}
  \big(b_{\mathrm{SM}}+\Delta b_{\mathrm{bridges}}\big)_i=(b_1,b_2,b_3)=(\tfrac{33}5,1,-3),
\end{equation}
the same triple that Theorem~\ref{thm:susy} obtains from the tower content. The two computations are
independent: one sums representations, the other sums bridges. \ProofWitnessBlock{bridges\_reproduce\_mssm, beta\_is\_mssm}

\end{corollary}

\begin{table}[htb]\centering\footnotesize
\caption{The $\beta$-coefficients as Standard Model plus ER=EPR bridges (exp38). Each bridge is
the inverted-statistics, same-representation companion of a SM field; $H_d$ is forced by the
higgsino anomaly (exp36). The bridges run as virtual degrees of freedom but are not asymptotic
states, so no superpartner appears in a detector.}\label{tab:susy}
\begin{tabular}{lccc}\hline
& $\delta b_3$ & $\delta b_2$ & $\delta b_1$\\\hline
$b_{\mathrm{SM}}$ & $-7$ & $-19/6$ & $41/10$\\\hline
gauginos (Weyl adjoint $\leftarrow$ gauge boson) & $2$ & $4/3$ & $0$\\
sfermions (scalar, same rep $\leftarrow$ fermion) & $2$ & $2$ & $2$\\
higgsino (Weyl $\leftarrow$ Higgs) & $0$ & $1/3$ & $1/5$\\
$H_d$ (anomaly-forced) & $0$ & $1/2$ & $3/10$\\\hline
$\Delta b_{\mathrm{bridges}}$ & $4$ & $25/6$ & $5/2$\\\hline
$b_{\mathrm{SM}}+\Delta b$ & $-3$ & $1$ & $33/5$\\\hline
\multicolumn{4}{c}{$(b_3,b_2,b_1)=(-3,\,1,\,33/5)=b_{\mathrm{MSSM}}$\ \checkmark}\\\hline
\end{tabular}\end{table}

\subsection{Third face: fluid}\label{ssec:fluid}
\begin{proposition}[Solitons, positroids, monopoles]\label{prop:fluid}
The regular KP solitons~\cite{KodamaWilliams} $u=2(\log\tau)_{xx}$ correspond exactly to points of $\Grpos$ (total
positivity $\Leftrightarrow\tau>0\Leftrightarrow$ regularity); their positroid profiles are the
condensate profiles $D(\sigma)-1=\varepsilon_0\golden^{\sigma}$, so the confining monopoles are
the bridges of \S\ref{ssec:susy}.  (Kodama--Williams~\cite{KodamaWilliams})
\ProofWitnessBlock{positroid\_from\_kp, bridge\_cocycle}
% TODO: no numerical test for positroid_from_kp, bridge_cocycle yet
\end{proposition}

\subsection{The gauge operators, AdS/CFT, and the mass gap}\label{ssec:operators}
The operator is the one \S\ref*{ssec:p-tower} defined (Remark~\ref{rmk:mass-operator}); here its
spectrum is read as masses and its gap is computed.

\begin{proposition}[Self-adjoint operators]\label{prop:operators}
The scale operator $H\lvert\sigma\rangle=m_0\golden^{\sigma}$ (\eqref{eq:mass-operator}; the golden
tower, not the Regge mass tower---Remark~\ref{rmk:mass-operator}) and the Yang--Mills Casimir $C_2(n)=n(n+3)/3$ are real diagonal, hence self-adjoint and bounded below,
with the vacuum at $0$. \ProofWitnessBlock{tower\_ham\_is\_hermitian, casimir\_op\_is\_hermitian, tower\_ham\_bounded\_below, casimir\_op\_bounded\_below}

\end{proposition}

\begin{proposition}[Discrete spectrum and the tower gap]\label{prop:spectrum}
The eigenvalues of $H$ increase without bound, $m_0\golden^{\sigma}\to\infty$ (since $\golden>1$
by \eqref{eq:phisq}), so
the resolvent is compact and the spectrum is discrete. The vacuum sits at $0$, every other level is
strictly positive, and the first excitation is separated by
\begin{equation}\label{eq:tower-gap}
  E_1-E_0=m_0>0 .
\end{equation}
This is a second gap, independent of the electric-Casimir gap of Theorem~\ref{thm:gap}: one is the
tower's own spacing, the other the colour energy. \ProofWitnessBlock{spectrum\_discrete, vacuum\_unique, mass\_gap, operators\_selfadjoint\_and\_isometry}

\end{proposition}
\begin{proof}
A real diagonal matrix equals its conjugate transpose, hence is Hermitian; its spectrum is its
diagonal, real and $\ge0$, with least value $0$.
\end{proof}

\begin{theorem}[The bulk Hamiltonian is the boundary modular generator]\label{thm:intertwine}
The bulk tower Hamiltonian and the tower entropy $S(\sigma)=\pi\golden^{\sigma}$ of
\eqref{eq:tower-modes} are exponentials of the same dilatation generator, at the same rate---the
regulator $R_K=\log\golden$ of Definition~\ref{def:regulator}:
\begin{equation}\label{eq:bulk-boundary-exp}
  H(\sigma)=m_0\,e^{\sigma R_K},
  \qquad
  S(\sigma)=\pi\,e^{\sigma R_K}.
\end{equation}
They differ only in the prefactor, so their ratio does not depend on the level:
\begin{equation}\label{eq:intertwine}
  \frac{\hat K(\sigma)}{H(\sigma)}=\frac{\pi}{m_0}
  \qquad\text{for every }\sigma,
  \qquad\text{equivalently}\qquad
  m_0\,\hat K(\sigma)=\pi\,H(\sigma).
\end{equation}
The bulk--boundary map of \eqref{eq:isometry}, an isometry acting level by level, therefore
intertwines them exactly. That $S$ is the boundary \emph{modular} generator---that the modular flow
is the tower's own scale flow---is established below, in \S\ref{ssec:zeta-gap}; with it,
\eqref{eq:intertwine} says the bulk Hamiltonian is that generator up to $\pi/m_0$.
(the GKPW dictionary~\cite{GKP98,Witten98})%
\ProofWitnessBlock{modular\_bulk\_ratio\_level\_independent, bulk\_boundary\_intertwine, ower\_e\_eq\_exp\_regulator, S\_tower\_eq\_exp\_regulator}%
\end{theorem}
\begin{proof}
$\golden^{\sigma}=e^{\sigma\log\golden}$ gives both displays in \eqref{eq:bulk-boundary-exp}. The
ratio is $\pi\golden^{\sigma}/(m_0\golden^{\sigma})=\pi/m_0$, and $\golden^{\sigma}>0$ cancels for
every $\sigma$.
\end{proof}

\begin{remark}[What this closes, and what it does not]\label{rmk:intertwine-scope}
This is the realisation \S\ref{ssec:adscft} deferred when it said that the microstate's
\emph{realisation as a self-adjoint operator} is \S\ref{sec:dof}. What the correspondence of
\S\ref{ssec:adscft} lacked was a concrete bulk Hamiltonian: self-adjoint, bounded below, with
discrete spectrum and a gap. Propositions~\ref{prop:operators} and~\ref{prop:spectrum} supply it,
and \eqref{eq:intertwine} says it is the tower entropy---and so, by \S\ref{ssec:zeta-gap}, the
boundary modular generator---up to one constant, which is why the correspondence holds \emph{at
each level} and not at a preferred one. Had the bulk
grown in any base other than $\golden$, the ratio would drift with $\sigma$ and no level-wise
isometry could intertwine them.

What is \emph{not} proved here is the GKPW dictionary itself; it enters cited~\cite{GKP98,Witten98}
(\cite{Maldacena98,GKP98,Witten98}), nor the identification of the modular flow with the
static-patch boost, which rests on Bisognano--Wichmann and is taken up in \S\ref{ssec:zeta-gap}.
\end{remark}

\begin{proposition}[The local gauge field (W4)]\label{prop:localfield}
The potential $A_\mu$ valued in $\su(3)=A_2$ (Proposition~\ref{prop:a2}) gives the local field
$F=dA+A\wedge A$ and the action $S=\tfrac1{4}\int\operatorname{tr}F^2$; the Wilson lattice measure
$e^{-S_W}$ on the $A_2$ lattice---the Eisenstein lattice $\Z[\omega]$ in which the eigenvalue
triangle sits (\S\ref{ssec:construction})---is gauge-invariant and reflection-positive, and generates
$T=e^{-aH}$. \ProofWitnessBlock{colour\_jacobi}
\end{proposition}

\begin{table}[htb]\centering\footnotesize
\caption{Derived spectra of the operators (first values).}\label{tab:spectra}
\begin{tabular}{ll}\hline
object & spectrum\\\hline
mode count $N(\sigma)=\lfloor\pi\golden^\sigma\rfloor$ & $3,\,5,\,8,\,13,\,21,\dots$\\
mass $H$ & $0,\,m_0,\,m_0\golden,\,m_0\golden^2,\dots$\\
Casimir $C_2(n)=n(n+3)/3$ & $0,\,\tfrac43,\,\tfrac{10}3,\,6,\dots$\\
transfer $T=e^{-aH}$ & $1,\,e^{-am_0},\,e^{-am_0\golden},\dots$\\\hline
\end{tabular}\end{table}

\begin{theorem}[The mass gap]\label{thm:gap}
The Kogut--Susskind transfer operator~\cite{KogutSusskind} built from the Wilson measure~\cite{Wilson} has a positive gap from the electric Casimir,
\begin{equation}\label{eq:gap}
  \Delta=\tfrac23 g^2>0. 
\end{equation}
\ProofWitnessBlock{strong\_coupling\_gap}
\end{theorem}
\begin{proof}
At strong coupling the transfer operator is diagonal in the electric-flux basis; the vacuum is the
trivial representation ($C_2=0$) and the first excited state carries the fundamental flux, with
electric energy $\tfrac23 g^2$ from $C_2$ of the fundamental. The difference is \eqref{eq:gap},
positive for $g^2>0$.
\end{proof}

\subsection{Reflection positivity and the OS reconstruction}\label{ssec:os}
A self-adjoint, bounded-below, gapped operator reconstructs a quantum field theory~\cite{OsterwalderSchrader}; the reflection
that makes the reconstruction positive is the one \S\ref*{ssec:hinge} already gave.

\begin{figure}[htb]\centering
\begin{tikzpicture}[font=\footnotesize,>={Stealth[length=4pt]}]
\node (H) at (0,0) {$H$};\node (T) at (3,0) {$T$};\node (O) at (6.4,0) {OS recon.};\node (V) at (9.8,0) {4D boundary};
\draw[->] (H)--node[above]{$e^{-aH}$} (T);
\draw[->] (T)--node[above]{$\Theta$} (O);
\draw[->] (O)--node[above]{$V$} (V);
\node[gray,below=0.5mm of H]{\scriptsize spectrum/masses};
\node[gray,below=0.5mm of T,align=center]{\scriptsize euclidean evol.\\[-2pt]RP\,+\,clustering};
\node[gray,below=0.5mm of O]{\scriptsize Wightman QFT};
\node[gray,below=0.5mm of V]{\scriptsize gauge theory};
\end{tikzpicture}
\caption{How the operators work together, all anchored on $\lvert\Om\rvert=\tfrac12$ (the vacuum of
$H$, the fixed point of $\Theta$, the norm $V$ preserves); $C_2$ explains the gap that $H$ carries.}
\label{fig:pipeline}
\end{figure}

\begin{proposition}[Reflection positivity]\label{prop:rp}
The transfer operator $T=e^{-aH}$ is positive, with eigenvalues $e^{-a m_0\golden^{\sigma}}>0$, so
\begin{equation}\label{eq:rp}
  \langle\Theta f,f\rangle=\langle f,Tf\rangle=\sum_{\sigma}|c_\sigma|^2\,e^{-a m_0\golden^{\sigma}}\ge0,
\end{equation}
where $\Theta$ is the antiunitary CPT/Galois involution of \S\ref*{ssec:hinge} ($\Theta^2=1$,
$\lVert\Theta z\rVert=\lVert z\rVert$), $F=e^{-(a/2)H}f$ is the state placed at half the
Euclidean separation from the reflection plane, and $f=\sum_\sigma c_\sigma\lvert\sigma\rangle$ is
arbitrary.
\ProofWitnessBlock{theta\_pairing\_eq\_transfer, reflection\_positive\_general, half\_prop\_sq, theta\_involutive, theta\_preserves\_modulus\_C, theta\_commutes\_real\_diagonal, transfer\_positive, cpt\_involution, cpt\_preserves\_modulus}

\end{proposition}
\begin{proof}
The content of \eqref{eq:rp} is the \emph{first} equality, not the inequality, and it is now
proved rather than described. Three facts do it. First, $\Theta$ is an antilinear involution of
modulus one that commutes with a real diagonal Hamiltonian: $\Theta^2=1$,
$\lVert\Theta z\rVert=\lVert z\rVert$, and $\Theta(Ez)=E\,\Theta z$ for $E$ real. Since $\Theta$
exchanges the two Rindler wedges and the wedge generator is the same on both, it fixes $H$, hence
fixes the tower basis and acts on a state by conjugating its coefficients. Second, the reflection
places $\Theta F$ at Euclidean distance $a/2$ on one side of the plane and $F$ at $a/2$ on the
other, and the two half-propagators compose,
\begin{equation}\label{eq:half-prop}
  e^{-(a/2)E_\sigma}\cdot e^{-(a/2)E_\sigma}=e^{-aE_\sigma}=T_\sigma ,
\end{equation}
which is where the separation being the KMS period $a=\beta=2\pi/H_{dS}$
(\S\ref*{ssec:accum}, W9) enters: it is what makes the total separation exactly $a$, and hence the
composite exactly $T$. Third, summing over the tower with $F=e^{-(a/2)H}f$ gives
$\langle\Theta F,F\rangle=\sum_\sigma c_\sigma^{2}T_\sigma=\langle f,Tf\rangle$ for \emph{every}
$f$, not term by term. The inequality then follows because each
$T_\sigma=e^{-am_0\golden^{\sigma}}>0$, so the sum is non-negative for every $f$ and every
truncation; positivity of $T$ is $H$ self-adjoint (Proposition~\ref{prop:operators}) with real
spectrum. $\Theta$ itself is \S\ref*{ssec:hinge}, W6.
\end{proof}

\begin{proposition}[Reflection positivity as a property of the measure]\label{prop:rp-measure}
The diagonal statement \eqref{eq:rp} is the transfer-operator form. The measure-theoretic form
holds as well, and for every linear functional supported in the half-space, not term by term.
For the one-dimensional free scalar of mass $m>0$, with covariance
$C(x,y)=e^{-m|x-y|}/(2\sinh m)$, points $x_i>0$ and arbitrary coefficients $c_i$,
\begin{equation}\label{eq:rp-measure}
  \sum_{i,j}c_ic_j\,C(-x_i,x_j)
  =\frac{1}{2\sinh m}\Bigl(\sum_i c_i e^{-m x_i}\Bigr)^{2}\;\ge\;0 .
\end{equation}
Equivalently the reflected covariance matrix is $(2\sinh m)^{-1}vv^{\mathsf T}$ with
$v_i=e^{-m x_i}$: a positive multiple of a rank-one Gram matrix, hence positive semidefinite.
\ProofWitnessBlock{refl\_cov\_factors, refl\_quad\_form\_eq\_square, reflection\_positive\_measure}
\end{proposition}
\begin{proof}
On the half-space $|-x_i-x_j|=x_i+x_j$, so $C(-x_i,x_j)=e^{-mx_i}e^{-mx_j}/(2\sinh m)$ and the
double sum factors into the square of a single sum. The factorisation is the whole content: the
reflection turns a difference of coordinates into a sum, and a covariance that was a function of
the difference becomes a product of one factor per point. The test discriminates: for a covariance
that is not of this form---an oscillating $C(x,y)=\cos 2(x-y)$, say---the reflected matrix has a
negative eigenvalue and reflection positivity fails, so \eqref{eq:rp-measure} is not automatic.
\end{proof}

\begin{theorem}[Osterwalder--Schrader reconstruction]\label{thm:os}
From $H$ the transfer operator $T=e^{-aH}$ yields, by \eqref{eq:rp}, a GNS Hilbert space with a
unique vacuum $\lvert\Om\rangle$ (top eigenvalue $1$, isolated by the gap $e^{-a m_0}<1$) and
exponential clustering $\langle A T^{n}B\rangle_c\sim e^{-a m_0 n}$. The von Neumann types of
\S\ref*{ssec:noselect} (I discrete, II modular, III$_1$ local) supply the remaining
Osterwalder--Schrader axioms. \ProofWitnessBlock{os\_reconstruction\_explicit, transfer\_spectral\_data, clustering\_exponential, vacuum\_transfer\_eq\_one, transfer\_gap\_below\_vacuum}

\end{theorem}
\begin{proof}
The sesquilinear form \eqref{eq:rp} is positive semidefinite; its GNS quotient and completion give
the Hilbert space, with the vacuum the eigenvector of $T$ of eigenvalue $1$ ($H\lvert\Om\rangle=0$).
Uniqueness and isolation are the gap $e^{-a m_0}<1$; clustering is the ratio $(e^{-a m_0})^{n}\to0$.
Reflection positivity is W6; locality is the type-III$_1$ non-factorisation
(\S\ref*{ssec:hinge}, Prop.~\ref*{prop:obs-nofirewall}); the discrete spectrum and the modular flow
are types I and II (\S\ref*{ssec:noselect}).
\end{proof}

\begin{corollary}[Cyclicity of the vacuum (W6)]\label{cor:cyclic}
The vacuum $\lvert\Om\rangle$ is cyclic: the tower states $\lvert\sigma\rangle$ are generated from
it by the field operators, so their span is dense in the GNS space.
\ProofWitnessBlock{os\_reconstruction\_explicit, tower\_ham\_is\_hermitian, tower\_ham\_bounded\_below}
\end{corollary}

\begin{remark}[Analyticity (OS0)]\label{rmk:os0}
The Euclidean--Lorentzian passage is the internal Wick rotation $(it)^2=-t^2$, geometric (carried
by the eigenvalue metric), not an external continuation; with the UV-finite worldsheet partition
(\S\ref*{ssec:p-tower}) it is the analyticity/temperedness axiom. \ProofWitnessBlock{internal\_wick}
\end{remark}

\begin{theorem}[The bulk--boundary isometry]\label{thm:isometry}
The bulk--boundary map $V$, $V^{\dagger}V=\mathbf1$---its norm core proved in Lemma~\ref{lem:isometry}, its holographic content \eqref{eq:isometry} of \S\ref{ssec:adscft}---together with the GKPW dictionary~\cite{Maldacena98,GKP98,Witten98}
$Z_{\mathrm{grav}}[\phi|_\partial]=\langle e^{\int_\partial\phi_0\mathcal O}\rangle$ and its isometric (quantum error-correcting) refinement~\cite{AlmheiriDongHarlow}, identifies the OS-reconstructed theory with the four-dimensional boundary
gauge theory. \ProofWitnessBlock{bulk\_boundary\_isometry}
\end{theorem}
\begin{proof}
$V^{\dagger}V=\mathbf1$ is Lemma~\ref{lem:isometry} (with \eqref{eq:isometry} of \S\ref{ssec:adscft}); an isometry preserves the inner product, so
bulk and boundary correlators coincide, which is the content of \eqref{eq:isometry}. The GKPW
equality then reads the boundary correlators as the reconstructed ones.
\end{proof}

\subsection{Confining spectrum, S-duality, and the surviving gap}\label{ssec:zeta-gap}
The spectrum the reconstruction produces is the tower already identified with $\zeta$, and the
duality that protects its gap is the reflection already proved in \S2.

\begin{theorem}[The confining spectrum is $\zeta$]\label{thm:confining}
The confining glueball spectrum is the closed charged bridges on the tower: the Regge tower
$=$ Euler product $=\zeta$ (Proposition~\ref{prop:veneziano}), the charge carried by the bridges
(\S\ref{ssec:susy}).
\ProofWitnessBlock{regge\_tower\_is\_euler\_product, regge\_dirichlet\_eq\_zeta}
\NumericalWitnessBlock{prop:veneziano}

\end{theorem}
\begin{proof}
A glueball is a closed flux tube; the flux tube is a closed bridge (\S\ref{ssec:fluid}); its
excitations lie on the Regge tower, whose Dirichlet series is $\zeta$
(Proposition~\ref{prop:veneziano}); the bridge carries the colour (\S\ref{ssec:susy}).
\end{proof}

\begin{proposition}[Magnetic condensate and the colour gap]\label{prop:condensate}
Apply the electric--magnetic duality of Theorem~\ref{thm:duality} to the holographic
superconductor already carried by the framework (\S\ref{ssec:accum}, in the de~Sitter setting of
\S\ref{ssec:observer}). Dirac quantisation
$q\,q_m=2\pi$ sends the electric condensate $V$ to a magnetic one, whose monopoles---the ER=EPR
bridges of \S\ref{ssec:susy}, in their fluid reading as the positroid profiles
$D(\sigma)-1=\varepsilon_0\golden^{\sigma}$ (Proposition~\ref{prop:fluid})---condense. Monopole
condensation confines colour--electric flux into a tube (dual Meissner), giving a string tension and
with it the gap
\begin{equation}\label{eq:string-tension}
  \sigma_{\mathrm{tension}}=q_m^{2}V>0,\qquad \Delta_{\mathrm{colour}}=\sqrt{\sigma_{\mathrm{tension}}}>0 .
\end{equation}
At the self-dual point $\tau=i$ the electric and magnetic gaps coincide, so the gap is invariant
under the duality that relates strong and weak coupling. ; the dual-superconductor mechanism \NumericalWitness{eq:string-tension}
is \cite{KogutSusskind} ('t~Hooft--Mandelstam).
\end{proposition}
\begin{proposition}[The tension is welded to the tower]\label{prop:tension-weld}
The condensate profile and the per-level energy are the same function,
$D(\sigma)-1=\varepsilon(\sigma)=\epsz\golden^{\sigma}$ (Proposition~\ref{prop:fluid},
\S\ref{ssec:accum}), while the potential runs the other way: the de~Sitter swampland condition
\eqref{eq:dS-swampland}, $|\partial_\sigma V|/V=\ln\golden$, integrates to
$V(\sigma)=\epsz\golden^{-\sigma}$\,\ProofWitness{swampland\_has\_deriv\_at}
The two are therefore conjugate,
\begin{equation}\label{eq:condensate-conjugate}
  V(\sigma)\,\big(D(\sigma)-1\big)=\epsz^{2}\qquad\text{for every }\sigma ,
\end{equation}
the same pattern as the throat/weld pair \eqref{eq:conjugate-pair}. Feeding this into
\eqref{eq:string-tension} with Dirac quantisation $q\,q_m=2\pi$ and the certainty relation
$\epsz\Mpcf=\pi$ gives a quantity invariant along the tower,
\begin{equation}\label{eq:tension-weld}
  \sigma_{\mathrm{tension}}(\sigma)\;S(\sigma)=q_m^{2}\,\epsz\,\pi=\frac{4\pi^{4}}{q^{2}\Mpcf},
\end{equation}
which is the gauge-side counterpart of the time--scale weld $S(\sigma)\tau_F(\sigma)=\pi\Mpcf$ of
\eqref{eq:obs-weld}: the string tension is not a free scale but is fixed level by level by the
saturation entropy, and $\Delta_{\mathrm{colour}}=\sqrt{\sigma_{\mathrm{tension}}}\propto\golden^{-\sigma/2}$
weakens towards the ultraviolet, as asymptotic freedom requires. Against the integer count
$\Nmodes=\lfloor S\rfloor$ the product carries the floor ripple, the same capacity/occupancy
distinction as \eqref{eq:obs-identity}. This identity is not a second entropic derivation running
parallel to $\delta Q=T\delta S$: on the positive object amplitude, fluid, degree of freedom and
state are one datum (Theorem~\ref{thm:faces}), so \eqref{eq:tension-weld} is an identity \emph{within}
that object rather than a bridge between two, and the dynamical link between the sectors is the
double copy $\text{graviton}=\mathrm{YM}^{2}$ of Theorem~\ref{thm:graviton}, with the field equation
itself obtained in \S\ref{app:einstein}. What \eqref{eq:tension-weld} adds is that the tension is
fixed level by level rather than free.
\ProofWitnessBlock{swampland\_has\_deriv\_at, eps0\_Mpcf\_pi, tension\_welded\_to\_tower, tension\_weld\_value}
\end{proposition}

\begin{proposition}[The two gaps are faces of the scale operator]\label{prop:gap-faces}
The scale operator $H\lvert\sigma\rangle=m_0\golden^{\sigma}$ (Proposition~\ref{prop:operators}) climbs the golden tower, and its saturation entropy is that spectrum up to the unit $\pi$, $S(\sigma)=\pi\golden^{\sigma}$. The welded string tension \eqref{eq:tension-weld} is the conjugate of that spectrum: $\sigma_{\mathrm{tension}}(\sigma)\propto\golden^{-\sigma}$ runs against the operator's $\golden^{\sigma}$, and the colour gap $\Delta_{\mathrm{colour}}=\sqrt{\sigma_{\mathrm{tension}}}\propto\golden^{-\sigma/2}$ descends the same tower the operator climbs. Hence the tower's own gap $E_1-E_0=m_0$ of \eqref{eq:tower-gap} and the electric--Casimir gap $\Delta=\tfrac23 g^2$ of Theorem~\ref{thm:gap} are two readings of the one operator's tower: one its spacing, the other the reciprocal-root of its conjugate.
\ProofWitnessBlock{S\_is\_operator\_spectrum, tension\_is\_conjugate\_spectrum, tension\_welded\_to\_tower, mass\_gap}
\end{proposition}

\begin{theorem}[The modular generator is the Landau--Lifshitz boundary charge]\label{thm:modular-LL}
The static patch is thermal at the Gibbons--Hawking temperature, and its modular flow is the tracial scale flow $\sigma\mapsto\sigma+t$ of the tower (\S\ref{ssec:accum})---a pure dilation $S(\sigma+t)=\golden^{t}S(\sigma)$, hence the static-patch boost. Because the ETS is a \emph{single} spacetime with tower levels, not a family of vacua, Bisognano--Wichmann is here not a black box but the tower's own scale group: the modular flow \emph{is} the scale flow, already proved a group preserving $\lvert\Om\rvert=\tfrac12$ (\S\ref{ssec:accum}). The modular Hamiltonian is therefore $\hat K_{\mathrm{mod}}=2\pi H_\xi/\kappa$, and with the horizon charge $H_\xi=\kappa A/(8\pi\GN)$ of Theorem~\ref{thm:LL-energy},
\[
  \hat K_{\mathrm{mod}}=\frac{A}{4\GN}=S(\sigma)=\pi\golden^{\sigma}.
\]
Matching the horizon entropy to the tower fixes $H(\sigma)=\sqrt2\,\golden^{-\sigma/2}$: the horizon descends the tower as the colour gap of Proposition~\ref{prop:gap-faces}, another reading of $\golden^{-\sigma/2}$. The crossed product turning the type~III algebra into type~II~\cite{CLPW22} is generated by $\hat K_{\mathrm{mod}}+q$, with $q$ the observer's energy; its gravitational part is the Landau--Lifshitz charge. The generator that makes the observer's algebra tracial is the tower entropy.
for the proved core---the modular flow is the tower's scale group, $\hat K_{\mathrm{mod}}=S(\sigma)$, and the descent $H(\sigma)\propto\golden^{-\sigma/2}$; Bisognano--Wichmann~\cite{BisognanoWichmann75} and the CLPW crossed product enter as hypotheses. \ProofWitnessBlock{modular\_generator\_is\_entropy, modular\_generator\_is\_LL\_over\_T, modular\_generator\_is\_tower, scale\_flow\_is\_dilation, hubble\_descends\_tower}
\end{theorem}

\begin{theorem}[The colour scale and running are derived from $M$]\label{thm:colour-from-M}
The lepton scale $M$ and the framework scale $\Mpcf$ coincide, both fixed by $\varepsilon_0(\cdot)=\pi$ \eqref{eq:obs-threshold}; hence the welded-tension invariant $4\pi^4/(q^2\Mpcf)$ of \eqref{eq:tension-weld} is set by the same $M$ that sets the charged-lepton masses, and the colour gap runs as fixed by the tower (Proposition~\ref{prop:gap-faces}). The colour scale and running are therefore determined by $M$, not free; the quark spectrum and CKM mixing are not fixed by this and remain open.
\ProofWitnessBlock{M\_eq\_Mpcf, colour\_scale\_from\_M, colour\_running\_fixed}
\end{theorem}

\begin{remark}[The two faces of $M$: dynamical and placement]\label{rmk:M-two-faces}
The scale $M=\Mpcf$ of Theorem~\ref{thm:colour-from-M} is the \emph{dynamical} face, exact by the certainty relation. It has a \emph{placement} face in the tower: with the proton at $\sigma=5$, five worldsheet transitions of action $\pi$ above the electron at $\sigma=0$, $m_p/m_e=\lvert S_3\rvert\,\pi^5=6\pi^5=1836.12$ (experiment $1836.15$), whence $m_p/3=312.8$~MeV. Two of the three factors are formal---$\pi=5\arccos(\golden/2)$ (Proposition~\ref{prop:pi-bridge}) and $\lvert S_3\rvert=3!=6$ (Lemma~\ref{lem:s3-orders})---while the assignment of the proton to $\sigma=5$ and the comparison of $6\pi^5$ with the measured ratio (to $1.9\times10^{-5}$) are empirical: a confrontation with data, not a theorem.
(data comparison only) \ProofWitnessBlock{M\_eq\_Mpcf}
\end{remark}

\begin{theorem}[One object, four registers]\label{thm:one-object}
Every scale in this framework is the single modulus $\lvert\Om\rvert=\mu_3=\tfrac12$ read in a different register. The hinge is that the certainty relation is that modulus: the half-turn $\pi$ that fixes the scale is $2\pi\mu_3$, so $\varepsilon_0\Mpcf=2\pi\mu_3$~\eqref{eq:certainty}. From this single identity the scale $M=\Mpcf=2\pi\mu_3/\varepsilon_0$ (Theorem~\ref{thm:colour-from-M}), the welded string tension carrying $4\pi^4/(q^2\Mpcf)$ \eqref{eq:tension-weld}, the colour scale and running that descend from it, and the proton placement $m_p/m_e=6\pi^5$ of the bridge all reduce to $\mu_3$: the phase of \S\ref{sec:methods}, the observer threshold~\eqref{eq:obs-threshold}, the welded tension, and the tower placement are one object in four registers.
\ProofWitnessBlock{scale\_is\_modulus, one\_object}
\end{theorem}

\begin{figure}[htb]\centering
\resizebox{\textwidth}{!}{%
\begin{tikzpicture}[
  font=\footnotesize, >={Stealth[length=4pt]},
  nd/.style={inner sep=2pt, outer sep=1pt},
  ap/.style={draw, rounded corners=2pt, fill=black!5, inner sep=3pt},
  every path/.style={->, draw=black!72, line width=0.45pt}]
  % raiz
  \node[ap] (om)  at (0,3.0) {$\lvert\Om\rvert=\mu_3=\tfrac12$};
  \node[nd] (cy)  at (0,1.6) {$\varepsilon_0\Mpcf=2\pi\mu_3=\pi$};
  % operadores / gap
  \node[nd] (H)   at (3.6,4.2) {$H\lvert\sigma\rangle=m_0\golden^\sigma$};
  \node[nd] (S)   at (3.6,3.0) {$S(\sigma)=\pi\golden^\sigma$};
  \node[nd] (gp)  at (3.6,1.8) {$E_1-E_0=m_0$};
  \node[nd] (cg)  at (3.6,0.6) {$\Delta=\tfrac23 g^2$};
  % tension soldada
  \node[nd] (tw)  at (7.4,3.6) {$\sigma_{\mathrm{tension}}S=\tfrac{4\pi^4}{q^2\Mpcf}$};
  \node[ap] (dc)  at (7.4,2.2) {$\Delta_{\mathrm{colour}}\propto\golden^{-\sigma/2}$};
  % escala / color / placement
  \node[nd] (Mm)  at (11.2,4.0) {$M=\Mpcf$};
  \node[ap] (col) at (11.2,2.6) {colour scale $+$ running};
  \node[nd] (p5)  at (11.2,1.2) {$m_p/m_e=6\pi^5$};
  % cierre
  \node[ap] (one) at (14.4,2.6) {one object};
  % flechas (las que saltan un nodo rodean por un lado)
  \draw (om)--(cy); \draw (cy)--(H); \draw (cy)--(S);
  \draw (H) to[bend right=28] (gp);
  \draw (S) to[bend left=28]  (cg);
  \draw (S)--(tw); \draw (cg)--(tw); \draw (tw)--(dc);
  \draw (cy) to[bend right=22] (Mm);
  \draw (Mm)--(col); \draw (dc)--(col);
  \draw (Mm) to[bend left=26] (p5);
  \draw (col)--(one); \draw (p5)--(one);
  \draw (dc) to[bend right=14] (one);
  \draw (om) to[bend right=38] (one);
\end{tikzpicture}}
\caption{Commutative spine of \S\ref{sec:dof} (operators, welded tension, colour): the same
microstate $\Om$, $\lvert\Om\rvert=\tfrac12$, propagates through the scale operator and its gap,
the welded tension, the colour scale and running, and the proton placement $6\pi^5$---every arrow
uses only earlier nodes, and all converge on the one modulus (Theorem~\ref{thm:one-object}).
Mathematical, string and gauge layers coincide on one diagram; shaded nodes are apices where
independent chains meet the same value.}
\label{fig:spine-dof}
\end{figure}

\begin{remark}[Which mechanism sets the glueball masses]\label{rmk:glueball-check}
Two distinct statements must be kept apart. Theorem~\ref{thm:confining} says which arithmetic object
\emph{indexes} the confining spectrum---the Regge tower, hence $\zeta$. It does not set mass values:
the poles of \eqref{eq:veneziano} carry the \emph{open} tower, with leading spin $J=n-1$, whereas a
glueball is a \emph{closed} flux tube. The masses come instead from
Proposition~\ref{prop:condensate}: the scale is the string tension,
$\Delta_{\mathrm{colour}}=\sqrt{\sigma_{\mathrm{tension}}}$, with the closed states obtained from the
open ones by the double copy of Theorem~\ref{thm:graviton}. Comparing $\sqrt n$ from the open tower
directly against lattice glueball ratios would ignore both the spin assignment $J\le n-1$ that
\S\ref{ssec:p-tower} imposes and the confining mechanism that actually fixes the scale. (the
spin-assignment check is \texttt{exp43b}; the condensate mechanism \texttt{exp47}).
\end{remark}

\begin{theorem}[S-duality is the functional equation]\label{thm:duality}
The framework's spectral partition is the completed zeta---the product over all places of
Theorem~\ref{thm:places}, its finite leg the Euler product that \eqref{eq:euler-product} identifies
with the Regge tower and its Archimedean leg the self-dual Gaussian of
Proposition~\ref{prop:selfdual-gaussian}---so the electric--magnetic duality
$s\mapsto1-s$---which the framework derives rather than imports, as the Galois conjugation
$\golden\leftrightarrow-1/\golden$ of the golden field, the two arithmetic conjugates being the
strong and weak coupling (\S\ref{ssec:web})---is the functional equation $\xi(s)=\xi(1-s)$ (Theorem~\ref{thm:funct-eq}), whose
unique fixed point $s=\tfrac12$ is the modulus $\lvert\Om\rvert$ (Remark~\ref{rmk:half-selfdual});
$\zeta$ itself is produced by $\pi$ (even values) and $\golden$ (odd values, via $\chi_5$;
\S\ref*{ssec:zeta}; also \cite{F1,Elvang2026}). The even values are Theorem~\ref{thm:even-zeta}, the odd ones
Theorem~\ref{thm:zeta-odd}, and the character Proposition~\ref{prop:pentagon-chi5}. \ProofWitnessBlock{s\_duality\_exact, functional\_equation\_fixed\_at\_half, places\_assembly}

\end{theorem}

\begin{corollary}[The gap survives the continuum]\label{cor:survive}
Along the one-loop asymptotic-freedom trajectory, on which the coupling runs as
$1/g^{2}(a)=b_0\log\!\big(1/(a\Lambda)\big)$, the transmuted scale is
\begin{equation}\label{eq:transmutation}
  \Lambda_{\mathrm{QCD}}(a)=a^{-1}e^{-1/(b_0g^{2}(a))}
  =a^{-1}e^{-\log(1/(a\Lambda))}=a^{-1}(a\Lambda)=\Lambda ,
\end{equation}
\emph{exactly, for every lattice spacing}. The physical gap is therefore independent of the cutoff,
$\Delta_{\mathrm{phys}}=\Lambda>0$. This is a consistency check on the one-loop parametrisation, not
the argument: the reason no continuum limit of the lattice kind arises here is structural and is given
in Remark~\ref{rmk:residual}. The scale
is fixed by $\varepsilon_0\Mpcf=\pi$ (\S\ref*{ssec:construction}) and the flow reaches the
ultraviolet fixed point $c(n)\to2$ (\S\ref*{ssec:accum}).
\ProofWitnessBlock{gap\_independent\_of\_cutoff, eps0\_Mpcf\_pi, lambda\_qcd\_eq\_lambda, lambda\_qcd\_pos}
\NumericalWitnessBlock{eq:transmutation}
\end{corollary}

\begin{remark}[Where the residual actually sits]\label{rmk:residual}
It is worth being exact about what is and is not settled here, because the pieces are stronger than
a blanket appeal to duality would suggest.

The strong-coupling gap \eqref{eq:gap} is the standard rigorous lattice result---Wilson measure,
Osterwalder--Seiler reflection positivity, Kogut--Susskind transfer operator---and does not depend on
any holographic model. \ProofWitnessBlock{strong\_coupling\_gap}

The duality is not borrowed from $\mathcal N=4$. The framework sits at \emph{its own} self-dual point
$\tau=i$, which is derived rather than chosen (\S\ref{ssec:construction}), and at a self-dual fixed
point electric and magnetic coincide by construction (Proposition~\ref{prop:condensate}); no
Montonen--Olive exactness is inherited. The duality's content is Riemann's functional equation, a
theorem (Theorem~\ref{thm:funct-eq}, Theorem~\ref{thm:duality}). \ProofWitnessBlock{fixed\_line\_is\_face\_phi, functional\_equation\_fixed\_line, places\_assembly, s\_duality\_exact, s\_duality\_fixed\_point\_unique, theta\_from\_jacobi, functional\_equation\_fixed\_at\_half, places\_assembly, s\_duality\_exact}

The magnetic sector is not posited either: the monopoles \emph{are} the ER=EPR bridges, whose
groupoid laws and positroid profiles are proven (Proposition~\ref{prop:er-epr-id},
Proposition~\ref{prop:fluid}). \ProofWitnessBlock{er\_epr\_bridge, er\_epr\_state\_eq\_geometry, bridge\_cocycle, positroid\_from\_kp}

Finally, the continuum limit does not arise here in the form it takes for lattice gauge theory, and
the reason is structural rather than technical. In the lattice construction the spacing $a$ is an
\emph{artificial ultraviolet regulator} introduced to tame divergences, and the continuum limit is the
problem of removing it while keeping the physics finite. In the present construction there is no such
ultraviolet regulator to remove:

\begin{itemize}
\item the ultraviolet is already finite---the one-loop worldsheet partition function
      \eqref{eq:pcf-partition} converges, regulated globally by the modular partition on the torus
      and locally at the vertex (\S\ref{ssec:p-tower}); no divergence calls for a cutoff;
\item the theory sits at an ultraviolet fixed point, the asymptotic-safety corner of the duality web
      \eqref{eq:bridge-fixed}, with the spectral flow reaching $c(n)\to2$
      (\S\ref{ssec:accum});
\item and that fixed point is not tuned: Assumption~\ref{asm:bridge} (\eqref{eq:obs-fixedpoint}) makes it \emph{equivalent} to the
      algebraic identity $\epsz\Mpcf=\pi$, which \S\ref{ssec:construction} derives from the half-turn
      accumulated once around the fundamental domain.
\end{itemize}

The discreteness of the tower is therefore not a regularisation of something continuous; it is the
structure, and the framework is defined at its own fixed point rather than approached from a
cutoff. Together with \eqref{eq:transmutation}---which holds exactly at every spacing---this leaves
no residual continuum problem of the lattice kind. The projection fixing exactly three generations, previously stated
here as an open input, is closed in the companion bridge: the matter sits on the aperiodic PCF
interface, where by total irrationality of $\golden$ (Hof's criterion) the cut cannot produce a
periodic lattice, so the GSO summation---a periodic-torus construction---does not apply, and the
$\F_1$-descent resolves the three even spin structures into the three observer components
$T_0,T_1,T_2=P\text{-}C\text{-}F$ rather than summing them. The count three is therefore not a
gauge-sector input at all: it is the arity of \S\ref{ssec:arity}, forced there by minimal
non-paradoxical self-reference---distributed recurrence requires depth $k\ge2$
(Definition~\ref{def:dsr}), depth $2$ is minimal with unique positive root $\golden$
(Theorem~\ref{thm:fib-min}), and $\golden^{2}+\golden^{-2}=3$ fixes the arity
---the same three that gives the colour arity $\lfloor\pi\rfloor$ and the
three eigenvalues $\tfrac12\omega^{k}$ of Proposition~\ref{prop:pcf-norms}. What the bridge supplies
is not the number but the resolution of the spin structures onto it.
\ProofWitnessBlock{framework\_mass\_gap, gap\_survives, eps0\_Mpcf\_pi}

\end{remark}

\begin{figure}[!tbp]\centering\footnotesize
{\scriptsize$\begin{array}{lccccc}
\textbf{standard} & G\ \text{(compact simple)} & \xrightarrow{\ \text{construct 4D QFT}\ } & \text{Wightman/OS axioms} & \xrightarrow{\ \text{spectrum}\ } & \Delta>0\\[1mm]
 & \updownarrow & & \updownarrow & & \updownarrow\\[1mm]
\textbf{here} & \su(3)=A_2 & \xrightarrow[\ \text{KS};\,\Grpos\ ]{\ \lfloor\pi\golden^\sigma\rfloor\ } & \text{reconstruction from } H,\,T,\,\Theta & \xrightarrow[\ \text{IR confining}\ ]{\ \text{lowest mode }3\neq0\ } & \Delta=\tfrac23g^{2}>0
\end{array}$}
\caption{The ordering criterion for this section, read as a \emph{contrast} rather than a claim of
identification: the standard route to a gapped four-dimensional gauge theory (top) beside the route
taken here (bottom), column by column. Where the standard construction posits a compact simple group,
$\su(3)=A_2$ is derived from the eigenvalue triangle (Proposition~\ref{prop:a2}); where it builds a
four-dimensional theory by regularisation, the tower and the positive Grassmannian supply one
directly; where it imposes the Wightman/OS axioms, they are recovered from the framework's own
operators $H$, $T=e^{-aH}$ and $\Theta$ (\S\ref{ssec:os}); and where it asks for a spectral gap, the
electric Casimir gives $\Delta=\tfrac23g^{2}$ (Theorem~\ref{thm:gap}), cutoff-independent by
\eqref{eq:transmutation}. Each column is the same requirement met two ways, not a correspondence
awaiting proof.}\label{fig:claysquare}
\end{figure}

\subsection{Why this is a single demonstration}\label{ssec:spine5}

\begin{theorem}[The section's checkable core]\label{thm:core}
The following hold together:
\begin{equation}\label{eq:core}
\begin{aligned}
  &\golden^2=\golden+1\ \eqref{eq:phisq}, &&\sigma_{\mathrm{spec}}=\tfrac32\ (\S\ref{ssec:spectrum}), &&\text{kissing}_i=2\,\text{posroots}_i,\\
  &(b_1,b_2,b_3)=(\tfrac{33}5,1,-3), &&\textstyle\sum Y=0, &&b_{\mathrm{SM}}+\Delta b_{\mathrm{bridges}}=-3,\\
  &\text{and for every } m_0>0: &&E_1-E_0=m_0>0 . &&{}
\end{aligned}
\end{equation}
Each conjunct is established above---$\golden^2=\golden+1$ in \S\ref{ssec:construction}, the kissing
identity in Proposition~\ref{prop:ladder}, the $\beta$-coefficients in Proposition~\ref{prop:sm} and
the anomaly sum with them, the bridge decomposition in Corollary~\ref{cor:bridges-mssm}, and the
spectral gap in Proposition~\ref{prop:spectrum}---and collected here they are the arithmetic core of
the section: each can be checked on its own without the surrounding reading. \ProofWitnessBlock{entropy\_dof\_core, gauge\_side\_facts, transfer\_spectral\_data}

\end{theorem}

The path went from entropy to degrees of freedom (\S\ref{ssec:law}), placed them on the positive
Grassmannian (\S\ref{ssec:grass}) and its ladder (\S\ref{ssec:ladder}), read them as amplitude
(\S\ref{ssec:amp}), particle and graviton (\S\ref{ssec:particle}), bridges (\S\ref{ssec:susy}) and
fluid (\S\ref{ssec:fluid}), then defined the gauge operator (Proposition~\ref{prop:operators}, with the local field of
Proposition~\ref{prop:localfield}) whose gap (\S\ref{ssec:operators}, Theorem~\ref{thm:gap} and
Corollary~\ref{cor:survive}),
reflection positivity (Proposition~\ref{prop:rp}) and reconstruction (\S\ref{ssec:os},
Theorem~\ref{thm:os}, Corollary~\ref{cor:cyclic}) give a Yang--Mills theory whose
confining spectrum is $\zeta$ and whose S-duality is $\zeta$'s functional equation, self-dual at
$\tfrac12$ (\S\ref{ssec:zeta-gap}). These are not eleven results but one. The same $\zeta$ is the
amplitude (\S\ref{ssec:amp}), the tower (\S\ref{ssec:law}), and the duality
(\S\ref{ssec:zeta-gap}); the same $\tfrac12$ is the modulus (\S\ref{ssec:grass}, Proposition~\ref{prop:two-pos}), the area-law
constant (\S\ref{ssec:law}), the reflection fixed point and the self-dual line
(\S\ref{ssec:os}--\ref{ssec:zeta-gap}). The section closes where \S\ref*{sec:methods} opened: the
self-dual $\tfrac12$ of the functional equation is the fixed point $\mu=\tfrac12$ of
$x\mapsto1-x$. The half returns to the half, as in \S\ref*{sec:implications} the non-local gravity
returned to its own projection.

\begin{figure}[!htbp]\centering
\resizebox{0.86\textwidth}{!}{%
\begin{tikzpicture}[font=\scriptsize,>={Stealth[length=3pt]},
  nd/.style={inner sep=1.5pt,outer sep=1pt},
  zt/.style={inner sep=1.5pt,outer sep=1pt,draw=Dandelion,fill=Dandelion!18,rounded corners=1pt},
  om/.style={inner sep=1.5pt,outer sep=1pt,draw=RoyalBlue,fill=RoyalBlue!12,rounded corners=1pt},
  hdr/.style={font=\bfseries\footnotesize},
  every path/.style={->,draw=black!62,line width=0.4pt}]

% ===== phase headers =====
\foreach \x/\t in {0/A amplitudes,3/B dictionary,6/C ladder,9/D Bianconi,12/E gravity,15/F particles/SUSY,18.2/G Yang--Mills}
  \node[hdr] at (\x,6.6) {\t};

% ===== A (x=0) =====
\node[nd] (A1) at (0,5.6){A1 Gr$_{\ge0}$};
\node[nd] (A2) at (0,4.6){A2 KP$\leftrightarrow$Gr$_{\ge0}$};
\node[nd] (A4) at (0,3.6){A4 $A_4{=}\tfrac{\Gamma\Gamma}\Gamma$};
\node[zt] (A5) at (0,2.6){A5 $\sum n^{-s}{=}\zeta$};
\node[nd] (A6) at (0,1.4){A6 $\golden{=}2\cos\tfrac\pi5$};
% ===== B (x=3) =====
\node[nd] (B1) at (3,5.2){B1 $\sigma{=}\tfrac32$};
\node[om] (B2) at (3,4.0){B2 $\lvert\Om\rvert{=}\tfrac12$};
\node[nd] (B3) at (3,2.8){B3 $\alpha'B{\to}\tfrac1s{+}\tfrac1t$};
\node[nd] (B5) at (3,1.6){B5 $\pi/\golden$ split};
% ===== C (x=6) =====
\node[nd] (C1) at (6,4.8){C1 $\tfrac12\omega^k{=}A_2$};
\node[nd] (C3) at (6,3.4){C3 $A_2{\subset}D_4{\subset}E_6{\subset}E_8$};
\node[nd] (C4) at (6,2.0){C4 dims$(2,4,6,8)$};
% ===== D (x=9) =====
\node[nd] (D1) at (9,5.2){D1 $\tfrac{\rm QFI}4{=}M_{\rm B}$};
\node[om] (D3) at (9,4.0){D3 $S{=}kA,k{=}\tfrac12$};
\node[zt] (D6) at (9,2.8){D6 DOF${\sim}\golden^\sigma$};
\node[nd] (D8) at (9,1.6){D8 4 faces $=1$};
% ===== E (x=12) =====
\node[nd] (E1) at (12,5.0){E1 $\square h{=}8\pi G T$};
\node[nd] (E5) at (12,3.6){E5 $R_{\mu\nu}{=}3H^2g$};
\node[nd] (E6) at (12,2.2){E6 flat${+}$curved dS};
% ===== F (x=15) =====
\node[nd] (F1) at (15,5.2){F1 $b{=}(\tfrac{33}5,1,{-}3)$};
\node[nd] (F2) at (15,4.0){F2 $b_{\rm MSSM}{=}b_{\rm SM}{+}$br};
\node[nd] (F5) at (15,2.8){F5 graviton${=}$YM$^2$};
\node[nd] (F4) at (15,1.6){F4 $\sin^2\theta_W|_{\rm GUT}{=}\tfrac38$};
% ===== G (x=18.2) =====
\node[nd] (G1) at (18.2,5.6){G1 $H\lvert\sigma\rangle{=}m_0\golden^\sigma$};
\node[nd] (G3) at (18.2,4.6){G3 $\Delta{=}\tfrac23g^2$};
\node[nd] (G5) at (18.2,3.6){G5 RP $\langle\theta f,f\rangle{\ge}0$};
\node[om] (G7) at (18.2,2.6){G7 $V^\dagger V{=}\mathbf1$};
\node[zt] (G9) at (18.2,1.6){G9 $\xi(s){=}\xi(1{-}s)$};
\node[nd] (G11) at (18.2,0.5){G11 $\Delta_{\rm phys}{=}\Lambda{>}0$};

% ===== arc A→B→C→D→E→F→G (representative edges) =====
\draw (A1)--(A2); \draw (A4)--(A5); \draw (A6)--(B2); \draw (A5)--(B1);
\draw (A4)--(B3); \draw (A2)--(B5);
\draw (B1)--(C1); \draw (B2)--(C1); \draw (C1)--(C3); \draw (C3)--(C4);
\draw (C1)--(D1); \draw (B2)--(D3); \draw (C3)--(D6); \draw (D1)--(D8); \draw (D6)--(D8);
\draw (D3)--(E1); \draw (D6)--(E5); \draw (E5)--(E6); \draw (D8)--(E1);
\draw (C3)--(F1); \draw (F1)--(F2); \draw (E1)--(F5); \draw (F1)--(F4);
\draw (F2)--(G1); \draw (F5)--(G1); \draw (D6)--(G1);
\draw (G1)--(G3); \draw (G1)--(G5); \draw (G5)--(G7); \draw (G3)--(G11);
\draw (G7)--(G9); \draw (G9)--(G11);

% ===== the two threading invariants (highlighted) =====
\draw[Dandelion,line width=1pt,densely dashed] (A5) to[out=-20,in=180] (D6);
\draw[Dandelion,line width=1pt,densely dashed] (D6) to[out=-10,in=170] (G9);
\draw[RoyalBlue,line width=1pt,densely dashed] (B2) to[out=-25,in=175] (D3);
\draw[RoyalBlue,line width=1pt,densely dashed] (D3) to[out=-15,in=185] (G7);
% closure: G9 self-dual 1/2 back to |Ω|=1/2
\draw[red!55,densely dotted,line width=0.7pt] (G9) to[out=-90,in=-90,looseness=0.35]
  node[below,font=\scriptsize]{self-dual $\tfrac12{=}\lvert\Om\rvert$} (B2);

\begin{scope}[on background layer]
\node[fit=(A1)(A6)(G1)(G11)(D1),draw=black!12,rounded corners,inner sep=6mm]{};\end{scope}
\node[font=\scriptsize,below=1mm] at (current bounding box.south)
  {\textcolor{Dandelion}{$\blacksquare$}\,$\zeta$-thread (A5$\to$D6$\to$G9)\quad\textcolor{RoyalBlue}{$\blacksquare$}\,$\lvert\Om\rvert{=}\tfrac12$-thread (B2$\to$D3$\to$G7)\quad\textcolor{red}{$\cdots$}\,closure (self-dual $\tfrac12$)};
\end{tikzpicture}
}
\caption{Commutative diagram of \S\ref{sec:dof}: its demonstrative skeleton across the seven movements (A amplitudes, B dictionary, C ladder, D Bianconi, E gravity, F particles/SUSY, G Yang--Mills), threaded by $\zeta$ (gold) and $\lvert\Om\rvert=\tfrac12$ (blue), closing on the self-dual $\tfrac12$.}\label{fig:spine5}\end{figure}

\FloatBarrier
